\documentclass[12pt,aspectratio=1610]{beamer}
\geometry{paperwidth=20cm,paperheight=12.5cm}
\usepackage[italian,provide=*]{babel}
\usepackage{lmodern,fontspec}
\setsansfont{NimbusSans-Regular.otf}[Path=sorgenti/font/,BoldFont=NimbusSans-Bold.otf,ItalicFont=NimbusSans-Italic.otf]
\setmainfont{NimbusSans-Regular.otf}[Path=sorgenti/font/,BoldFont=NimbusSans-Bold.otf,ItalicFont=NimbusSans-Italic.otf]
\setmonofont{lmmono10-regular.otf}[Path=sorgenti/font/,Scale=.88]
\usefonttheme{professionalfonts}
\usepackage{amsmath,amssymb,booktabs,array,tabularx,graphicx,tikz,listings,bookmark}
\usetikzlibrary{calc,arrows.meta,positioning}
\definecolor{Paper}{HTML}{F8F6F1}
\definecolor{Ink}{HTML}{253344}\definecolor{Navy}{HTML}{15283E}
\definecolor{Quantum}{HTML}{5146A5}\definecolor{Classical}{HTML}{087C74}
\definecolor{Measure}{HTML}{A05B10}\definecolor{Rose}{HTML}{9B4564}
\definecolor{Muted}{HTML}{657180}\definecolor{Rule}{HTML}{D6DAD9}
\definecolor{Wash}{HTML}{EAE7F6}\definecolor{Night}{HTML}{112236}
\definecolor{Lilac}{HTML}{ABB5FF}\definecolor{Mint}{HTML}{72D8C8}\definecolor{Gold}{HTML}{EEC184}
\setbeamercolor{background canvas}{bg=Paper}\setbeamercolor{normal text}{fg=Ink}
\setbeamertemplate{navigation symbols}{}\setbeamertemplate{headline}{}\setbeamertemplate{footline}{}
\setbeamersize{text margin left=0pt,text margin right=0pt}\setbeameroption{hide notes}
\newsavebox{\BodyBox}\newsavebox{\SideBox}\newsavebox{\RecallBox}\newsavebox{\TitleBox}
\newcommand{\MeasureBox}[2]{\typeout{GEOBOX:#1:\the\wd#2:\the\ht#2:\the\dp#2}}
\lstset{language=R,basicstyle=\ttfamily\fontsize{10.2}{13.4}\selectfont,
columns=fullflexible,keepspaces=true,showstringspaces=false,breaklines=true,breakatwhitespace=false,
commentstyle=\color{Muted},keywordstyle=\color{Classical},stringstyle=\color{Quantum},
aboveskip=0pt,belowskip=0pt,xleftmargin=0pt,xrightmargin=0pt,tabsize=2}
\hypersetup{colorlinks=true,linkcolor=Classical,urlcolor=Classical,pdfpagemode=UseOutlines,pdfstartview=Fit,bookmarksopen=true,bookmarksopenlevel=0}
\bookmarksetup{open,openlevel=0,depth=2}
\tikzset{arr/.style={-{Stealth[length=2mm]},line width=.8pt,draw=Classical}}
\hypersetup{pdftitle={Campionamento, stima e incertezza},pdfauthor={},pdfsubject={Corso di statistica e data science per Geologia, primo anno; modulo 05}}
\usetikzlibrary{decorations.pathreplacing,fit,backgrounds,matrix}\hypersetup{pdftitle={M11 · Statistica spaziale},pdfsubject={Revisione integrale R6}}\begin{document}
\begin{frame}[plain]\hypertarget{slide1}{}\begin{tikzpicture}[remember picture,overlay]\fill[Night] (current page.south west) rectangle (current page.north east);\draw[Mint,line width=2pt] ([xshift=1.05cm,yshift=-1.35cm]current page.north west)--([xshift=3cm,yshift=-1.35cm]current page.north west);\node[anchor=north west,inner sep=0] at ([xshift=1.05cm,yshift=-1.7cm]current page.north west){\color{Mint}\bfseries\fontsize{11}{14}\selectfont M11 · REVISIONE R6};
\node[anchor=north west,inner sep=0] at ([xshift=1.05cm,yshift=-2.65cm]current page.north west){\begin{minipage}{15.8cm}\raggedright\color{white}\bfseries\fontsize{31}{38}\selectfont Statistica spaziale\end{minipage}};
\node[anchor=north west,inner sep=0] at ([xshift=1.05cm,yshift=-6.4cm]current page.north west){\begin{minipage}{15.8cm}\raggedright\color{Lilac}\fontsize{15}{21}\selectfont Coordinate, associazione e previsione nel territorio\end{minipage}};
\node[anchor=north west,inner sep=0] at ([xshift=1.05cm,yshift=-8.6cm]current page.north west){\begin{minipage}{16.5cm}\color{white}\fontsize{11}{15}\selectfont 32 slide · teoria, derivazioni, figure e applicazioni in R.\par\smallskip Provenienza dei dati e sorgenti nei materiali allegati.\end{minipage}};
\node[anchor=south east,inner sep=0] at ([xshift=18.95cm,yshift=-11.35cm]current page.north west){\color{Lilac!45}\bfseries\fontsize{64}{68}\selectfont M11};
\end{tikzpicture}\end{frame}
\begin{frame}[plain]\hypertarget{slide2}{}\bookmark[level=0,page=2]{Indice}\begin{tikzpicture}[remember picture,overlay]\node[anchor=north west,inner sep=0] at ([xshift=0.95cm,yshift=-0.65cm]current page.north west){\color{Navy}\bfseries\fontsize{22}{27}\selectfont Percorso del modulo};
\draw[Classical,line width=1.3pt] ([xshift=.95cm,yshift=-1.68cm]current page.north west)--([xshift=2.25cm,yshift=-1.68cm]current page.north west);\node[anchor=north west,inner sep=0] at ([xshift=0.95cm,yshift=-2.23cm]current page.north west){\color{Classical}\bfseries\fontsize{13}{16}\selectfont 3};
\node[anchor=north west,inner sep=0] at ([xshift=2.0cm,yshift=-2.23cm]current page.north west){\begin{minipage}{7.25cm}\raggedright\color{Navy}\bfseries\fontsize{12.4}{15.5}\selectfont\hyperlink{slide3}{Dato, posizione e supporto}\end{minipage}};
\node[anchor=north west,inner sep=0] at ([xshift=0.95cm,yshift=-3.5300000000000002cm]current page.north west){\color{Classical}\bfseries\fontsize{13}{16}\selectfont 6};
\node[anchor=north west,inner sep=0] at ([xshift=2.0cm,yshift=-3.5300000000000002cm]current page.north west){\begin{minipage}{7.25cm}\raggedright\color{Navy}\bfseries\fontsize{12.4}{15.5}\selectfont\hyperlink{slide6}{Coordinate e geometria}\end{minipage}};
\node[anchor=north west,inner sep=0] at ([xshift=0.95cm,yshift=-4.83cm]current page.north west){\color{Classical}\bfseries\fontsize{13}{16}\selectfont 10};
\node[anchor=north west,inner sep=0] at ([xshift=2.0cm,yshift=-4.83cm]current page.north west){\begin{minipage}{7.25cm}\raggedright\color{Navy}\bfseries\fontsize{12.4}{15.5}\selectfont\hyperlink{slide10}{Carte, appartenenze e aggregazioni}\end{minipage}};
\node[anchor=north west,inner sep=0] at ([xshift=0.95cm,yshift=-6.130000000000001cm]current page.north west){\color{Classical}\bfseries\fontsize{13}{16}\selectfont 13};
\node[anchor=north west,inner sep=0] at ([xshift=2.0cm,yshift=-6.130000000000001cm]current page.north west){\begin{minipage}{7.25cm}\raggedright\color{Navy}\bfseries\fontsize{12.4}{15.5}\selectfont\hyperlink{slide13}{Vicinati e matrici dei pesi}\end{minipage}};
\node[anchor=north west,inner sep=0] at ([xshift=0.95cm,yshift=-7.43cm]current page.north west){\color{Classical}\bfseries\fontsize{13}{16}\selectfont 16};
\node[anchor=north west,inner sep=0] at ([xshift=2.0cm,yshift=-7.43cm]current page.north west){\begin{minipage}{7.25cm}\raggedright\color{Navy}\bfseries\fontsize{12.4}{15.5}\selectfont\hyperlink{slide16}{Associazione spaziale e riferimento inferenziale}\end{minipage}};
\node[anchor=north west,inner sep=0] at ([xshift=10.299999999999999cm,yshift=-2.23cm]current page.north west){\color{Classical}\bfseries\fontsize{13}{16}\selectfont 21};
\node[anchor=north west,inner sep=0] at ([xshift=11.35cm,yshift=-2.23cm]current page.north west){\begin{minipage}{7.25cm}\raggedright\color{Navy}\bfseries\fontsize{12.4}{15.5}\selectfont\hyperlink{slide21}{Semivarianza e struttura a distanza}\end{minipage}};
\node[anchor=north west,inner sep=0] at ([xshift=10.299999999999999cm,yshift=-3.5300000000000002cm]current page.north west){\color{Classical}\bfseries\fontsize{13}{16}\selectfont 23};
\node[anchor=north west,inner sep=0] at ([xshift=11.35cm,yshift=-3.5300000000000002cm]current page.north west){\begin{minipage}{7.25cm}\raggedright\color{Navy}\bfseries\fontsize{12.4}{15.5}\selectfont\hyperlink{slide23}{Copertura e unità della sintesi}\end{minipage}};
\node[anchor=north west,inner sep=0] at ([xshift=10.299999999999999cm,yshift=-4.83cm]current page.north west){\color{Classical}\bfseries\fontsize{13}{16}\selectfont 25};
\node[anchor=north west,inner sep=0] at ([xshift=11.35cm,yshift=-4.83cm]current page.north west){\begin{minipage}{7.25cm}\raggedright\color{Navy}\bfseries\fontsize{12.4}{15.5}\selectfont\hyperlink{slide25}{Interpolazione e verifica geografica}\end{minipage}};
\node[anchor=north west,inner sep=0] at ([xshift=10.299999999999999cm,yshift=-6.130000000000001cm]current page.north west){\color{Classical}\bfseries\fontsize{13}{16}\selectfont 31};
\node[anchor=north west,inner sep=0] at ([xshift=11.35cm,yshift=-6.130000000000001cm]current page.north west){\begin{minipage}{7.25cm}\raggedright\color{Navy}\bfseries\fontsize{12.4}{15.5}\selectfont\hyperlink{slide31}{Consegna e studio autonomo}\end{minipage}};
\node[anchor=north west,inner sep=0] at ([xshift=0.95cm,yshift=-10.55cm]current page.north west){\begin{minipage}{18cm}\color{Muted}\fontsize{10}{13}\selectfont I numeri indicano le pagine iniziali; i titoli sono cliccabili. Codice, figure e verifiche sono inclusi nel progetto.\end{minipage}};
\end{tikzpicture}\end{frame}
\begin{lrbox}{\TitleBox}\begin{minipage}[t]{18.2cm}\raggedright\sloppy\color{Ink}\fontsize{20.7}{24}\selectfont \bfseries\color{Navy}La disposizione spaziale delle osservazioni\end{minipage}\end{lrbox}
\begin{lrbox}{\BodyBox}\begin{minipage}[t]{10.25cm}\raggedright\sloppy\color{Ink}\fontsize{12.4}{16.4}\selectfont Un’indagine geologica associa a ciascuna misura una posizione e un supporto, che può essere un volume di carota, un affioramento o un’area campionata. Due insiemi aventi gli stessi valori possono descrivere configurazioni territoriali differenti, poiché la distribuzione marginale ignora la loro disposizione. Il capitolo utilizza sessantaquattro siti sintetici su un reticolo locale di passo 250 metri, nei quali la concentrazione è espressa in unità arbitrarie. La configurazione consente di esaminare separatamente rappresentazione cartografica, associazione e previsione.\par\vspace{.28cm}{\centering\fontsize{13}{17}\selectfont\(\displaystyle\begin{gathered}\mathcal D=\{(s_i,z_i,A_i):i=1,\ldots,n\}\end{gathered}\)\par}\end{minipage}\end{lrbox}
\begin{lrbox}{\SideBox}\begin{minipage}[t]{7.45cm}\raggedright\sloppy\color{Ink}\fontsize{12.2}{16.1}\selectfont \includegraphics[width=7.45cm]{figure/M11/F01.pdf}\par\vspace{.16cm}{\fontsize{9.5}{12.6}\selectfont $s_i$: posizione; $z_i$: misura; $A_i$: supporto. Reticolo 8 × 8; concentrazione in unità arbitrarie (u.a.).}\end{minipage}\end{lrbox}
\begin{lrbox}{\RecallBox}\begin{minipage}[t]{16.45cm}\raggedright\sloppy\color{Ink}\fontsize{10.1}{13.2}\selectfont Posizione, valore e supporto appartengono alla definizione dell’osservazione.\end{minipage}\end{lrbox}
\MeasureBox{3-TitleBox}{\TitleBox}
\MeasureBox{3-BodyBox}{\BodyBox}
\MeasureBox{3-SideBox}{\SideBox}
\MeasureBox{3-RecallBox}{\RecallBox}
\begin{frame}[plain]\hypertarget{slide3}{}
\bookmark[level=0,page=3]{Dato, posizione e supporto}
\bookmark[level=1,page=3]{La disposizione spaziale delle osservazioni}
\begin{tikzpicture}[remember picture,overlay]\node[anchor=north east,inner sep=0] at ([xshift=19.05cm,yshift=-0.16cm]current page.north west){\color{Classical}\fontsize{8}{10}\selectfont\hyperlink{slide2}{Indice}};
\node[anchor=north west,inner sep=0] at ([xshift=0.95cm,yshift=-0.62cm]current page.north west){\usebox{\TitleBox}};
\draw[Classical,line width=1.3pt] ([xshift=.95cm,yshift=-1.68cm]current page.north west)--([xshift=2.25cm,yshift=-1.68cm]current page.north west);\node[anchor=north west,inner sep=0] at ([xshift=0.95cm,yshift=-2.12cm]current page.north west){\usebox{\BodyBox}};
\node[anchor=north west,inner sep=0] at ([xshift=11.6cm,yshift=-2.1cm]current page.north west){\color{Quantum}\bfseries\fontsize{11.5}{14}\selectfont Evidenza grafica};
\node[anchor=north west,inner sep=0] at ([xshift=11.6cm,yshift=-2.75cm]current page.north west){\usebox{\SideBox}};
\draw[Rule,line width=.5pt] ([xshift=.95cm,yshift=-10.43cm]current page.north west)--([xshift=19.05cm,yshift=-10.43cm]current page.north west);\node[anchor=north west,inner sep=0] at ([xshift=0.95cm,yshift=-10.66cm]current page.north west){\color{Classical}\bfseries\fontsize{9.4}{12}\selectfont Richiami};
\node[anchor=north west,inner sep=0] at ([xshift=2.6cm,yshift=-10.64cm]current page.north west){\usebox{\RecallBox}};
\draw[Rule,line width=.35pt] ([xshift=.95cm,yshift=-11.92cm]current page.north west)--([xshift=19.05cm,yshift=-11.92cm]current page.north west);\node[anchor=north west,inner sep=0] at ([xshift=0.95cm,yshift=-12.06cm]current page.north west){\color{Muted}\fontsize{7.7}{9.3}\selectfont M11 · Statistica spaziale};
\node[anchor=north east,inner sep=0] at ([xshift=17.65cm,yshift=-12.06cm]current page.north west){\color{Classical}\fontsize{7.3}{9.1}\selectfont };
\node[anchor=north east,inner sep=0] at ([xshift=19.05cm,yshift=-12.06cm]current page.north west){\color{Muted}\fontsize{7.7}{9.3}\selectfont 3/32};
\end{tikzpicture}\end{frame}
\begin{lrbox}{\TitleBox}\begin{minipage}[t]{18.2cm}\raggedright\sloppy\color{Ink}\fontsize{20.7}{24}\selectfont \bfseries\color{Navy}Distribuzione marginale e configurazione geografica\end{minipage}\end{lrbox}
\begin{lrbox}{\BodyBox}\begin{minipage}[t]{10.25cm}\raggedright\sloppy\color{Ink}\fontsize{12.4}{16.4}\selectfont Permutando le concentrazioni fra i siti si conservano media, varianza, quantili e istogramma, mentre cambiano le relazioni fra misure vicine. Nel reticolo di 64 siti a passo 250 m, la disposizione casuale modifica la continuità visiva senza alterare i valori misurati. Tale operazione diventerà un riferimento inferenziale soltanto dopo avere specificato l’ipotesi di scambiabilità. Una semplice permutazione grafica serve già a distinguere le proprietà dell’insieme dei valori dalle proprietà della loro organizzazione nello spazio.\par\vspace{.28cm}{\centering\fontsize{13}{17}\selectfont\(\displaystyle\begin{gathered}\{z_{\pi(i)}\}_{i=1}^n=\{z_i\}_{i=1}^n\quad\text{come multinsiemi}\end{gathered}\)\par}\end{minipage}\end{lrbox}
\begin{lrbox}{\SideBox}\begin{minipage}[t]{7.45cm}\raggedright\sloppy\color{Ink}\fontsize{12.2}{16.1}\selectfont \includegraphics[width=7.45cm]{figure/M11/F02.pdf}\par\vspace{.16cm}{\fontsize{9.5}{12.6}\selectfont $\pi$ riordina gli indici. Valori permutati; colori nella medesima scala di concentrazione, in u.a.}\end{minipage}\end{lrbox}
\begin{lrbox}{\RecallBox}\begin{minipage}[t]{16.45cm}\raggedright\sloppy\color{Ink}\fontsize{10.1}{13.2}\selectfont Una carta introduce informazione che un istogramma non conserva.\end{minipage}\end{lrbox}
\MeasureBox{4-TitleBox}{\TitleBox}
\MeasureBox{4-BodyBox}{\BodyBox}
\MeasureBox{4-SideBox}{\SideBox}
\MeasureBox{4-RecallBox}{\RecallBox}
\begin{frame}[plain]\hypertarget{slide4}{}
\bookmark[level=1,page=4]{Distribuzione marginale e configurazione geografica}
\begin{tikzpicture}[remember picture,overlay]\node[anchor=north east,inner sep=0] at ([xshift=19.05cm,yshift=-0.16cm]current page.north west){\color{Classical}\fontsize{8}{10}\selectfont\hyperlink{slide2}{Indice}};
\node[anchor=north west,inner sep=0] at ([xshift=0.95cm,yshift=-0.62cm]current page.north west){\usebox{\TitleBox}};
\draw[Classical,line width=1.3pt] ([xshift=.95cm,yshift=-1.68cm]current page.north west)--([xshift=2.25cm,yshift=-1.68cm]current page.north west);\node[anchor=north west,inner sep=0] at ([xshift=0.95cm,yshift=-2.12cm]current page.north west){\usebox{\BodyBox}};
\node[anchor=north west,inner sep=0] at ([xshift=11.6cm,yshift=-2.1cm]current page.north west){\color{Quantum}\bfseries\fontsize{11.5}{14}\selectfont Evidenza grafica};
\node[anchor=north west,inner sep=0] at ([xshift=11.6cm,yshift=-2.75cm]current page.north west){\usebox{\SideBox}};
\draw[Rule,line width=.5pt] ([xshift=.95cm,yshift=-10.43cm]current page.north west)--([xshift=19.05cm,yshift=-10.43cm]current page.north west);\node[anchor=north west,inner sep=0] at ([xshift=0.95cm,yshift=-10.66cm]current page.north west){\color{Classical}\bfseries\fontsize{9.4}{12}\selectfont Richiami};
\node[anchor=north west,inner sep=0] at ([xshift=2.6cm,yshift=-10.64cm]current page.north west){\usebox{\RecallBox}};
\draw[Rule,line width=.35pt] ([xshift=.95cm,yshift=-11.92cm]current page.north west)--([xshift=19.05cm,yshift=-11.92cm]current page.north west);\node[anchor=north west,inner sep=0] at ([xshift=0.95cm,yshift=-12.06cm]current page.north west){\color{Muted}\fontsize{7.7}{9.3}\selectfont M11 · Statistica spaziale};
\node[anchor=north east,inner sep=0] at ([xshift=17.65cm,yshift=-12.06cm]current page.north west){\color{Classical}\fontsize{7.3}{9.1}\selectfont };
\node[anchor=north east,inner sep=0] at ([xshift=19.05cm,yshift=-12.06cm]current page.north west){\color{Muted}\fontsize{7.7}{9.3}\selectfont 4/32};
\end{tikzpicture}\end{frame}
\begin{lrbox}{\TitleBox}\begin{minipage}[t]{18.2cm}\raggedright\sloppy\color{Ink}\fontsize{20.7}{24}\selectfont \bfseries\color{Navy}Il supporto determina la quantità osservata\end{minipage}\end{lrbox}
\begin{lrbox}{\BodyBox}\begin{minipage}[t]{10.25cm}\raggedright\sloppy\color{Ink}\fontsize{12.4}{16.4}\selectfont Se una misura rappresenta la media entro un’area, il suo confronto con una misura puntuale richiede un modello del cambiamento di supporto. L’aggregazione attenua variazioni locali e modifica la covarianza, anche quando la grandezza fisica e l’unità rimangono identiche. Una carta di valori medi regionali non permette quindi di ricostruire le differenze interne a ciascuna regione. Nel reticolo didattico i siti vengono trattati come osservazioni puntuali ideali, rendendo esplicita una semplificazione che, nelle indagini reali, deve essere giustificata dal protocollo.\par\vspace{.28cm}{\centering\fontsize{13}{17}\selectfont\(\displaystyle\begin{gathered}Z_A=\frac{1}{|A|}\int_A Z(s)\,\mathrm ds\end{gathered}\)\par}\end{minipage}\end{lrbox}
\begin{lrbox}{\SideBox}\begin{minipage}[t]{7.45cm}\raggedright\sloppy\color{Ink}\fontsize{12.2}{16.1}\selectfont \input{tikz/M11/D01.tex}\par\vspace{.16cm}{\fontsize{9.5}{12.6}\selectfont A: area di osservazione; $|A|$: sua superficie; Z(s): campo nella posizione s. $Z_A$ è la media areale.}\end{minipage}\end{lrbox}
\begin{lrbox}{\RecallBox}\begin{minipage}[t]{16.45cm}\raggedright\sloppy\color{Ink}\fontsize{10.1}{13.2}\selectfont L’unità statistica deve mantenere il medesimo significato lungo l’analisi.\end{minipage}\end{lrbox}
\MeasureBox{5-TitleBox}{\TitleBox}
\MeasureBox{5-BodyBox}{\BodyBox}
\MeasureBox{5-SideBox}{\SideBox}
\MeasureBox{5-RecallBox}{\RecallBox}
\begin{frame}[plain]\hypertarget{slide5}{}
\bookmark[level=1,page=5]{Il supporto determina la quantità osservata}
\begin{tikzpicture}[remember picture,overlay]\node[anchor=north east,inner sep=0] at ([xshift=19.05cm,yshift=-0.16cm]current page.north west){\color{Classical}\fontsize{8}{10}\selectfont\hyperlink{slide2}{Indice}};
\node[anchor=north west,inner sep=0] at ([xshift=0.95cm,yshift=-0.62cm]current page.north west){\usebox{\TitleBox}};
\draw[Classical,line width=1.3pt] ([xshift=.95cm,yshift=-1.68cm]current page.north west)--([xshift=2.25cm,yshift=-1.68cm]current page.north west);\node[anchor=north west,inner sep=0] at ([xshift=0.95cm,yshift=-2.12cm]current page.north west){\usebox{\BodyBox}};
\node[anchor=north west,inner sep=0] at ([xshift=11.6cm,yshift=-2.1cm]current page.north west){\color{Quantum}\bfseries\fontsize{11.5}{14}\selectfont Lettura del diagramma};
\node[anchor=north west,inner sep=0] at ([xshift=11.6cm,yshift=-2.75cm]current page.north west){\usebox{\SideBox}};
\draw[Rule,line width=.5pt] ([xshift=.95cm,yshift=-10.43cm]current page.north west)--([xshift=19.05cm,yshift=-10.43cm]current page.north west);\node[anchor=north west,inner sep=0] at ([xshift=0.95cm,yshift=-10.66cm]current page.north west){\color{Classical}\bfseries\fontsize{9.4}{12}\selectfont Richiami};
\node[anchor=north west,inner sep=0] at ([xshift=2.6cm,yshift=-10.64cm]current page.north west){\usebox{\RecallBox}};
\draw[Rule,line width=.35pt] ([xshift=.95cm,yshift=-11.92cm]current page.north west)--([xshift=19.05cm,yshift=-11.92cm]current page.north west);\node[anchor=north west,inner sep=0] at ([xshift=0.95cm,yshift=-12.06cm]current page.north west){\color{Muted}\fontsize{7.7}{9.3}\selectfont M11 · Statistica spaziale};
\node[anchor=north east,inner sep=0] at ([xshift=17.65cm,yshift=-12.06cm]current page.north west){\color{Classical}\fontsize{7.3}{9.1}\selectfont };
\node[anchor=north east,inner sep=0] at ([xshift=19.05cm,yshift=-12.06cm]current page.north west){\color{Muted}\fontsize{7.7}{9.3}\selectfont 5/32};
\end{tikzpicture}\end{frame}
\begin{lrbox}{\TitleBox}\begin{minipage}[t]{18.2cm}\raggedright\sloppy\color{Ink}\fontsize{20.7}{24}\selectfont \bfseries\color{Navy}Coordinate locali e riferimento geografico\end{minipage}\end{lrbox}
\begin{lrbox}{\BodyBox}\begin{minipage}[t]{10.25cm}\raggedright\sloppy\color{Ink}\fontsize{12.4}{16.4}\selectfont Le colonne x\_m e y\_m descrivono un piano locale con origine convenzionale, mentre lon e lat forniscono coordinate geografiche sintetiche prossime a 9° E e 45° N. La somiglianza delle estensioni non autorizza ad assegnare al piano locale un codice EPSG: il riferimento deve provenire dai metadati. Un controllo plausibile esamina identificativi, valori finiti, unità, ordine degli assi e collocazione attesa. Invertendo longitudine e latitudine si ottengono numeri globalmente ammissibili che collocano tuttavia i siti in un’altra regione.\par\vspace{.28cm}{\centering\fontsize{13}{17}\selectfont\(\displaystyle\begin{gathered}s_i=(x_i,y_i),\qquad [x_i]=[y_i]=\mathrm m\end{gathered}\)\par}\end{minipage}\end{lrbox}
\begin{lrbox}{\SideBox}\begin{minipage}[t]{7.45cm}\raggedright\sloppy\color{Ink}\fontsize{12.2}{16.1}\selectfont \input{tikz/M11/D02.tex}\end{minipage}\end{lrbox}
\begin{lrbox}{\RecallBox}\begin{minipage}[t]{16.45cm}\raggedright\sloppy\color{Ink}\fontsize{10.1}{13.2}\selectfont La validità numerica delle coordinate deve essere accompagnata da una verifica geografica.\end{minipage}\end{lrbox}
\MeasureBox{6-TitleBox}{\TitleBox}
\MeasureBox{6-BodyBox}{\BodyBox}
\MeasureBox{6-SideBox}{\SideBox}
\MeasureBox{6-RecallBox}{\RecallBox}
\begin{frame}[plain]\hypertarget{slide6}{}
\bookmark[level=0,page=6]{Coordinate e geometria}
\bookmark[level=1,page=6]{Coordinate locali e riferimento geografico}
\begin{tikzpicture}[remember picture,overlay]\node[anchor=north east,inner sep=0] at ([xshift=19.05cm,yshift=-0.16cm]current page.north west){\color{Classical}\fontsize{8}{10}\selectfont\hyperlink{slide2}{Indice}};
\node[anchor=north west,inner sep=0] at ([xshift=0.95cm,yshift=-0.62cm]current page.north west){\usebox{\TitleBox}};
\draw[Classical,line width=1.3pt] ([xshift=.95cm,yshift=-1.68cm]current page.north west)--([xshift=2.25cm,yshift=-1.68cm]current page.north west);\node[anchor=north west,inner sep=0] at ([xshift=0.95cm,yshift=-2.12cm]current page.north west){\usebox{\BodyBox}};
\node[anchor=north west,inner sep=0] at ([xshift=11.6cm,yshift=-2.1cm]current page.north west){\color{Quantum}\bfseries\fontsize{11.5}{14}\selectfont Lettura del diagramma};
\node[anchor=north west,inner sep=0] at ([xshift=11.6cm,yshift=-2.75cm]current page.north west){\usebox{\SideBox}};
\draw[Rule,line width=.5pt] ([xshift=.95cm,yshift=-10.43cm]current page.north west)--([xshift=19.05cm,yshift=-10.43cm]current page.north west);\node[anchor=north west,inner sep=0] at ([xshift=0.95cm,yshift=-10.66cm]current page.north west){\color{Classical}\bfseries\fontsize{9.4}{12}\selectfont Richiami};
\node[anchor=north west,inner sep=0] at ([xshift=2.6cm,yshift=-10.64cm]current page.north west){\usebox{\RecallBox}};
\draw[Rule,line width=.35pt] ([xshift=.95cm,yshift=-11.92cm]current page.north west)--([xshift=19.05cm,yshift=-11.92cm]current page.north west);\node[anchor=north west,inner sep=0] at ([xshift=0.95cm,yshift=-12.06cm]current page.north west){\color{Muted}\fontsize{7.7}{9.3}\selectfont M11 · Statistica spaziale};
\node[anchor=north east,inner sep=0] at ([xshift=17.65cm,yshift=-12.06cm]current page.north west){\color{Classical}\fontsize{7.3}{9.1}\selectfont };
\node[anchor=north east,inner sep=0] at ([xshift=19.05cm,yshift=-12.06cm]current page.north west){\color{Muted}\fontsize{7.7}{9.3}\selectfont 6/32};
\end{tikzpicture}\end{frame}
\begin{lrbox}{\TitleBox}\begin{minipage}[t]{18.2cm}\raggedright\sloppy\color{Ink}\fontsize{20.7}{24}\selectfont \bfseries\color{Navy}Assegnazione del CRS e trasformazione\end{minipage}\end{lrbox}
\begin{lrbox}{\BodyBox}\begin{minipage}[t]{10.25cm}\raggedright\sloppy\color{Ink}\fontsize{12.4}{16.4}\selectfont L’assegnazione del sistema di riferimento dichiara come interpretare numeri già presenti; la trasformazione calcola una nuova rappresentazione della medesima posizione. Il codice costruisce punti con ordine longitudine–latitudine esplicito e li trasforma in WGS 84 / UTM zona 32N. Il controllo di ritorno verifica la coerenza numerica dell’operazione entro una tolleranza prefissata. Tale controllo potrebbe riuscire anche partendo da un riferimento iniziale errato, poiché verifica la reversibilità della trasformazione scelta e non la correttezza dei metadati di partenza.\end{minipage}\end{lrbox}
\begin{lrbox}{\SideBox}\begin{minipage}[t]{7.45cm}\raggedright\sloppy\color{Ink}\fontsize{12.2}{16.1}\selectfont \par\vspace{.25cm}\lstinputlisting[language=R]{estratti/M11/riferimento.R}\end{minipage}\end{lrbox}
\begin{lrbox}{\RecallBox}\begin{minipage}[t]{16.45cm}\raggedright\sloppy\color{Ink}\fontsize{10.1}{13.2}\selectfont Il controllo di ritorno integra la verifica dei metadati, dai quali dipende il significato della trasformazione.\end{minipage}\end{lrbox}
\MeasureBox{7-TitleBox}{\TitleBox}
\MeasureBox{7-BodyBox}{\BodyBox}
\MeasureBox{7-SideBox}{\SideBox}
\MeasureBox{7-RecallBox}{\RecallBox}
\begin{frame}[plain]\hypertarget{slide7}{}
\bookmark[level=1,page=7]{Assegnazione del CRS e trasformazione}
\begin{tikzpicture}[remember picture,overlay]\node[anchor=north east,inner sep=0] at ([xshift=19.05cm,yshift=-0.16cm]current page.north west){\color{Classical}\fontsize{8}{10}\selectfont\hyperlink{slide2}{Indice}};
\node[anchor=north west,inner sep=0] at ([xshift=0.95cm,yshift=-0.62cm]current page.north west){\usebox{\TitleBox}};
\draw[Classical,line width=1.3pt] ([xshift=.95cm,yshift=-1.68cm]current page.north west)--([xshift=2.25cm,yshift=-1.68cm]current page.north west);\node[anchor=north west,inner sep=0] at ([xshift=0.95cm,yshift=-2.12cm]current page.north west){\usebox{\BodyBox}};
\node[anchor=north west,inner sep=0] at ([xshift=11.6cm,yshift=-2.1cm]current page.north west){\color{Quantum}\bfseries\fontsize{11.5}{14}\selectfont Implementazione in R};
\node[anchor=north west,inner sep=0] at ([xshift=11.6cm,yshift=-2.75cm]current page.north west){\usebox{\SideBox}};
\draw[Rule,line width=.5pt] ([xshift=.95cm,yshift=-10.43cm]current page.north west)--([xshift=19.05cm,yshift=-10.43cm]current page.north west);\node[anchor=north west,inner sep=0] at ([xshift=0.95cm,yshift=-10.66cm]current page.north west){\color{Classical}\bfseries\fontsize{9.4}{12}\selectfont Richiami};
\node[anchor=north west,inner sep=0] at ([xshift=2.6cm,yshift=-10.64cm]current page.north west){\usebox{\RecallBox}};
\draw[Rule,line width=.35pt] ([xshift=.95cm,yshift=-11.92cm]current page.north west)--([xshift=19.05cm,yshift=-11.92cm]current page.north west);\node[anchor=north west,inner sep=0] at ([xshift=0.95cm,yshift=-12.06cm]current page.north west){\color{Muted}\fontsize{7.7}{9.3}\selectfont M11 · Statistica spaziale};
\node[anchor=north east,inner sep=0] at ([xshift=17.65cm,yshift=-12.06cm]current page.north west){\color{Classical}\fontsize{7.3}{9.1}\selectfont sf: st\_transform};
\node[anchor=north east,inner sep=0] at ([xshift=19.05cm,yshift=-12.06cm]current page.north west){\color{Muted}\fontsize{7.7}{9.3}\selectfont 7/32};
\end{tikzpicture}\end{frame}
\begin{lrbox}{\TitleBox}\begin{minipage}[t]{18.2cm}\raggedright\sloppy\color{Ink}\fontsize{20.7}{24}\selectfont \bfseries\color{Navy}Distanze nel piano e sulla superficie terrestre\end{minipage}\end{lrbox}
\begin{lrbox}{\BodyBox}\begin{minipage}[t]{10.25cm}\raggedright\sloppy\color{Ink}\fontsize{12.4}{16.4}\selectfont Nel piano locale si impiega la distanza euclidea, la cui unità coincide con quella delle coordinate. Per coordinate angolari occorre invece scegliere una geometria sferica o ellissoidica, oppure una proiezione adeguata alla regione. La formula sferica usa angoli in radianti e un raggio terrestre dichiarato; un calcolo ellissoidico può restituire valori leggermente differenti. Nella matrice delle distanze, simmetria e diagonale nulla sono controlli necessari, mentre la loro soddisfazione non identifica la geometria più appropriata all’indagine.\par\vspace{.28cm}{\centering\fontsize{13}{17}\selectfont\(\displaystyle\begin{gathered}d_{ij}=\sqrt{(x_i-x_j)^2+(y_i-y_j)^2}\end{gathered}\)\par}\end{minipage}\end{lrbox}
\begin{lrbox}{\SideBox}\begin{minipage}[t]{7.45cm}\raggedright\sloppy\color{Ink}\fontsize{12.2}{16.1}\selectfont Nel reticolo, la diagonale fra i due angoli opposti vale 1750√2 m, circa 2475 m. La matrice D contiene 64 × 64 elementi; le distanze fra coppie distinte sono 2016 quando ogni coppia viene contata una sola volta.\end{minipage}\end{lrbox}
\begin{lrbox}{\RecallBox}\begin{minipage}[t]{16.45cm}\raggedright\sloppy\color{Ink}\fontsize{10.1}{13.2}\selectfont L’unità della distanza attraversa soglie, buffer, pesi e interpretazione dei risultati.\end{minipage}\end{lrbox}
\MeasureBox{8-TitleBox}{\TitleBox}
\MeasureBox{8-BodyBox}{\BodyBox}
\MeasureBox{8-SideBox}{\SideBox}
\MeasureBox{8-RecallBox}{\RecallBox}
\begin{frame}[plain]\hypertarget{slide8}{}
\bookmark[level=1,page=8]{Distanze nel piano e sulla superficie terrestre}
\begin{tikzpicture}[remember picture,overlay]\node[anchor=north east,inner sep=0] at ([xshift=19.05cm,yshift=-0.16cm]current page.north west){\color{Classical}\fontsize{8}{10}\selectfont\hyperlink{slide2}{Indice}};
\node[anchor=north west,inner sep=0] at ([xshift=0.95cm,yshift=-0.62cm]current page.north west){\usebox{\TitleBox}};
\draw[Classical,line width=1.3pt] ([xshift=.95cm,yshift=-1.68cm]current page.north west)--([xshift=2.25cm,yshift=-1.68cm]current page.north west);\node[anchor=north west,inner sep=0] at ([xshift=0.95cm,yshift=-2.12cm]current page.north west){\usebox{\BodyBox}};
\node[anchor=north west,inner sep=0] at ([xshift=11.6cm,yshift=-2.1cm]current page.north west){\color{Quantum}\bfseries\fontsize{11.5}{14}\selectfont Sviluppo e interpretazione};
\node[anchor=north west,inner sep=0] at ([xshift=11.6cm,yshift=-2.75cm]current page.north west){\usebox{\SideBox}};
\draw[Rule,line width=.5pt] ([xshift=.95cm,yshift=-10.43cm]current page.north west)--([xshift=19.05cm,yshift=-10.43cm]current page.north west);\node[anchor=north west,inner sep=0] at ([xshift=0.95cm,yshift=-10.66cm]current page.north west){\color{Classical}\bfseries\fontsize{9.4}{12}\selectfont Richiami};
\node[anchor=north west,inner sep=0] at ([xshift=2.6cm,yshift=-10.64cm]current page.north west){\usebox{\RecallBox}};
\draw[Rule,line width=.35pt] ([xshift=.95cm,yshift=-11.92cm]current page.north west)--([xshift=19.05cm,yshift=-11.92cm]current page.north west);\node[anchor=north west,inner sep=0] at ([xshift=0.95cm,yshift=-12.06cm]current page.north west){\color{Muted}\fontsize{7.7}{9.3}\selectfont M11 · Statistica spaziale};
\node[anchor=north east,inner sep=0] at ([xshift=17.65cm,yshift=-12.06cm]current page.north west){\color{Classical}\fontsize{7.3}{9.1}\selectfont };
\node[anchor=north east,inner sep=0] at ([xshift=19.05cm,yshift=-12.06cm]current page.north west){\color{Muted}\fontsize{7.7}{9.3}\selectfont 8/32};
\end{tikzpicture}\end{frame}
\begin{lrbox}{\TitleBox}\begin{minipage}[t]{18.2cm}\raggedright\sloppy\color{Ink}\fontsize{20.7}{24}\selectfont \bfseries\color{Navy}Distanza sferica e unità angolari\end{minipage}\end{lrbox}
\begin{lrbox}{\BodyBox}\begin{minipage}[t]{10.25cm}\raggedright\sloppy\color{Ink}\fontsize{12.4}{16.4}\selectfont La distanza lungo un arco di cerchio massimo si ricava da latitudini e longitudini espresse in radianti. Nella formula, a raccoglie la separazione angolare e R rappresenta il raggio della sfera adottata. Il laboratorio usa il modello sferico come confronto esplicito con la geometria piana; la precisione di un calcolo territoriale reale dipende anche dal datum, dall’estensione e dalla scelta fra sfera ed ellissoide. Le convenzioni geometriche devono accompagnare il risultato.\par\vspace{.28cm}{\centering\fontsize{13}{17}\selectfont\(\displaystyle\begin{gathered}a=\sin^2\frac{\Delta\varphi}{2}+\cos\varphi_i\cos\varphi_j\sin^2\frac{\Delta\lambda}{2}\\ d=2R\arcsin\sqrt a\end{gathered}\)\par}\end{minipage}\end{lrbox}
\begin{lrbox}{\SideBox}\begin{minipage}[t]{7.45cm}\raggedright\sloppy\color{Ink}\fontsize{12.2}{16.1}\selectfont $\varphi$: latitudine; $\lambda$: longitudine; $\Delta$ indica la differenza fra i siti i e j. Da gradi a radianti si moltiplica per $\pi/180$. R è il raggio sferico e determina l’unità della distanza. Inserire gradi nelle funzioni trigonometriche di R altera il significato del calcolo.\end{minipage}\end{lrbox}
\begin{lrbox}{\RecallBox}\begin{minipage}[t]{16.45cm}\raggedright\sloppy\color{Ink}\fontsize{10.1}{13.2}\selectfont La conversione angolare è parte del procedimento di calcolo e deve essere verificata.\end{minipage}\end{lrbox}
\MeasureBox{9-TitleBox}{\TitleBox}
\MeasureBox{9-BodyBox}{\BodyBox}
\MeasureBox{9-SideBox}{\SideBox}
\MeasureBox{9-RecallBox}{\RecallBox}
\begin{frame}[plain]\hypertarget{slide9}{}
\bookmark[level=1,page=9]{Distanza sferica e unità angolari}
\begin{tikzpicture}[remember picture,overlay]\node[anchor=north east,inner sep=0] at ([xshift=19.05cm,yshift=-0.16cm]current page.north west){\color{Classical}\fontsize{8}{10}\selectfont\hyperlink{slide2}{Indice}};
\node[anchor=north west,inner sep=0] at ([xshift=0.95cm,yshift=-0.62cm]current page.north west){\usebox{\TitleBox}};
\draw[Classical,line width=1.3pt] ([xshift=.95cm,yshift=-1.68cm]current page.north west)--([xshift=2.25cm,yshift=-1.68cm]current page.north west);\node[anchor=north west,inner sep=0] at ([xshift=0.95cm,yshift=-2.12cm]current page.north west){\usebox{\BodyBox}};
\node[anchor=north west,inner sep=0] at ([xshift=11.6cm,yshift=-2.1cm]current page.north west){\color{Quantum}\bfseries\fontsize{11.5}{14}\selectfont Sviluppo e interpretazione};
\node[anchor=north west,inner sep=0] at ([xshift=11.6cm,yshift=-2.75cm]current page.north west){\usebox{\SideBox}};
\draw[Rule,line width=.5pt] ([xshift=.95cm,yshift=-10.43cm]current page.north west)--([xshift=19.05cm,yshift=-10.43cm]current page.north west);\node[anchor=north west,inner sep=0] at ([xshift=0.95cm,yshift=-10.66cm]current page.north west){\color{Classical}\bfseries\fontsize{9.4}{12}\selectfont Richiami};
\node[anchor=north west,inner sep=0] at ([xshift=2.6cm,yshift=-10.64cm]current page.north west){\usebox{\RecallBox}};
\draw[Rule,line width=.35pt] ([xshift=.95cm,yshift=-11.92cm]current page.north west)--([xshift=19.05cm,yshift=-11.92cm]current page.north west);\node[anchor=north west,inner sep=0] at ([xshift=0.95cm,yshift=-12.06cm]current page.north west){\color{Muted}\fontsize{7.7}{9.3}\selectfont M11 · Statistica spaziale};
\node[anchor=north east,inner sep=0] at ([xshift=17.65cm,yshift=-12.06cm]current page.north west){\color{Classical}\fontsize{7.3}{9.1}\selectfont };
\node[anchor=north east,inner sep=0] at ([xshift=19.05cm,yshift=-12.06cm]current page.north west){\color{Muted}\fontsize{7.7}{9.3}\selectfont 9/32};
\end{tikzpicture}\end{frame}
\begin{lrbox}{\TitleBox}\begin{minipage}[t]{18.2cm}\raggedright\sloppy\color{Ink}\fontsize{20.7}{24}\selectfont \bfseries\color{Navy}Scala cromatica, geometria e copertura\end{minipage}\end{lrbox}
\begin{lrbox}{\BodyBox}\begin{minipage}[t]{10.25cm}\raggedright\sloppy\color{Ink}\fontsize{12.4}{16.4}\selectfont La carta di punti conserva il luogo nel quale la misura è stata acquisita, mentre una superficie continua richiede una regola di previsione fra i siti. Il rapporto unitario fra gli assi evita che una direzione venga deformata rispetto all’altra. Confronti fra carte richiedono la stessa scala cromatica e gli stessi limiti, affinché un colore mantenga il medesimo significato numerico. Anche i vuoti di campionamento devono rimanere riconoscibili, perché la densità dei simboli informa sulla copertura dell’indagine.\end{minipage}\end{lrbox}
\begin{lrbox}{\SideBox}\begin{minipage}[t]{7.45cm}\raggedright\sloppy\color{Ink}\fontsize{12.2}{16.1}\selectfont \includegraphics[width=7.45cm]{figure/M11/F01.pdf}\par\vspace{.16cm}{\fontsize{9.5}{12.6}\selectfont Coordinate in metri, proporzioni uguali; misure nei soli siti.}\end{minipage}\end{lrbox}
\begin{lrbox}{\RecallBox}\begin{minipage}[t]{16.45cm}\raggedright\sloppy\color{Ink}\fontsize{10.1}{13.2}\selectfont La grammatica cartografica deve conservare unità, geometria e copertura.\end{minipage}\end{lrbox}
\MeasureBox{10-TitleBox}{\TitleBox}
\MeasureBox{10-BodyBox}{\BodyBox}
\MeasureBox{10-SideBox}{\SideBox}
\MeasureBox{10-RecallBox}{\RecallBox}
\begin{frame}[plain]\hypertarget{slide10}{}
\bookmark[level=0,page=10]{Carte, appartenenze e aggregazioni}
\bookmark[level=1,page=10]{Scala cromatica, geometria e copertura}
\begin{tikzpicture}[remember picture,overlay]\node[anchor=north east,inner sep=0] at ([xshift=19.05cm,yshift=-0.16cm]current page.north west){\color{Classical}\fontsize{8}{10}\selectfont\hyperlink{slide2}{Indice}};
\node[anchor=north west,inner sep=0] at ([xshift=0.95cm,yshift=-0.62cm]current page.north west){\usebox{\TitleBox}};
\draw[Classical,line width=1.3pt] ([xshift=.95cm,yshift=-1.68cm]current page.north west)--([xshift=2.25cm,yshift=-1.68cm]current page.north west);\node[anchor=north west,inner sep=0] at ([xshift=0.95cm,yshift=-2.12cm]current page.north west){\usebox{\BodyBox}};
\node[anchor=north west,inner sep=0] at ([xshift=11.6cm,yshift=-2.1cm]current page.north west){\color{Quantum}\bfseries\fontsize{11.5}{14}\selectfont Evidenza grafica};
\node[anchor=north west,inner sep=0] at ([xshift=11.6cm,yshift=-2.75cm]current page.north west){\usebox{\SideBox}};
\draw[Rule,line width=.5pt] ([xshift=.95cm,yshift=-10.43cm]current page.north west)--([xshift=19.05cm,yshift=-10.43cm]current page.north west);\node[anchor=north west,inner sep=0] at ([xshift=0.95cm,yshift=-10.66cm]current page.north west){\color{Classical}\bfseries\fontsize{9.4}{12}\selectfont Richiami};
\node[anchor=north west,inner sep=0] at ([xshift=2.6cm,yshift=-10.64cm]current page.north west){\usebox{\RecallBox}};
\draw[Rule,line width=.35pt] ([xshift=.95cm,yshift=-11.92cm]current page.north west)--([xshift=19.05cm,yshift=-11.92cm]current page.north west);\node[anchor=north west,inner sep=0] at ([xshift=0.95cm,yshift=-12.06cm]current page.north west){\color{Muted}\fontsize{7.7}{9.3}\selectfont M11 · Statistica spaziale};
\node[anchor=north east,inner sep=0] at ([xshift=17.65cm,yshift=-12.06cm]current page.north west){\color{Classical}\fontsize{7.3}{9.1}\selectfont };
\node[anchor=north east,inner sep=0] at ([xshift=19.05cm,yshift=-12.06cm]current page.north west){\color{Muted}\fontsize{7.7}{9.3}\selectfont 10/32};
\end{tikzpicture}\end{frame}
\begin{lrbox}{\TitleBox}\begin{minipage}[t]{18.2cm}\raggedright\sloppy\color{Ink}\fontsize{20.7}{24}\selectfont \bfseries\color{Navy}Partizioni territoriali e appartenenze multiple\end{minipage}\end{lrbox}
\begin{lrbox}{\BodyBox}\begin{minipage}[t]{10.25cm}\raggedright\sloppy\color{Ink}\fontsize{12.4}{16.4}\selectfont La divisione Ovest–Est del laboratorio assegna all’Ovest i punti con x inferiore a 1000 metri e all’Est tutti gli altri, cosicché ogni identificativo compaia una sola volta. Due aree sovrapposte definiscono invece una relazione molti a molti: un punto può appartenere a entrambe. Con A definita da x minore o uguale a 1250 e B da x maggiore o uguale a 750, le sessantaquattro misure generano ottantotto relazioni. Una media calcolata sulla tabella delle relazioni attribuirebbe peso doppio ai ventiquattro punti comuni.\par\vspace{.28cm}{\centering\fontsize{13}{17}\selectfont\(\displaystyle\begin{gathered}|A|+|B|=|A\cup B|+|A\cap B|=64+24=88\end{gathered}\)\par}\end{minipage}\end{lrbox}
\begin{lrbox}{\SideBox}\begin{minipage}[t]{7.45cm}\raggedright\sloppy\color{Ink}\fontsize{12.2}{16.1}\selectfont \input{tikz/M11/D03.tex}\end{minipage}\end{lrbox}
\begin{lrbox}{\RecallBox}\begin{minipage}[t]{16.45cm}\raggedright\sloppy\color{Ink}\fontsize{10.1}{13.2}\selectfont Il denominatore dipende dall’unità della tabella: punti, aree o relazioni.\end{minipage}\end{lrbox}
\MeasureBox{11-TitleBox}{\TitleBox}
\MeasureBox{11-BodyBox}{\BodyBox}
\MeasureBox{11-SideBox}{\SideBox}
\MeasureBox{11-RecallBox}{\RecallBox}
\begin{frame}[plain]\hypertarget{slide11}{}
\bookmark[level=1,page=11]{Partizioni territoriali e appartenenze multiple}
\begin{tikzpicture}[remember picture,overlay]\node[anchor=north east,inner sep=0] at ([xshift=19.05cm,yshift=-0.16cm]current page.north west){\color{Classical}\fontsize{8}{10}\selectfont\hyperlink{slide2}{Indice}};
\node[anchor=north west,inner sep=0] at ([xshift=0.95cm,yshift=-0.62cm]current page.north west){\usebox{\TitleBox}};
\draw[Classical,line width=1.3pt] ([xshift=.95cm,yshift=-1.68cm]current page.north west)--([xshift=2.25cm,yshift=-1.68cm]current page.north west);\node[anchor=north west,inner sep=0] at ([xshift=0.95cm,yshift=-2.12cm]current page.north west){\usebox{\BodyBox}};
\node[anchor=north west,inner sep=0] at ([xshift=11.6cm,yshift=-2.1cm]current page.north west){\color{Quantum}\bfseries\fontsize{11.5}{14}\selectfont Lettura del diagramma};
\node[anchor=north west,inner sep=0] at ([xshift=11.6cm,yshift=-2.75cm]current page.north west){\usebox{\SideBox}};
\draw[Rule,line width=.5pt] ([xshift=.95cm,yshift=-10.43cm]current page.north west)--([xshift=19.05cm,yshift=-10.43cm]current page.north west);\node[anchor=north west,inner sep=0] at ([xshift=0.95cm,yshift=-10.66cm]current page.north west){\color{Classical}\bfseries\fontsize{9.4}{12}\selectfont Richiami};
\node[anchor=north west,inner sep=0] at ([xshift=2.6cm,yshift=-10.64cm]current page.north west){\usebox{\RecallBox}};
\draw[Rule,line width=.35pt] ([xshift=.95cm,yshift=-11.92cm]current page.north west)--([xshift=19.05cm,yshift=-11.92cm]current page.north west);\node[anchor=north west,inner sep=0] at ([xshift=0.95cm,yshift=-12.06cm]current page.north west){\color{Muted}\fontsize{7.7}{9.3}\selectfont M11 · Statistica spaziale};
\node[anchor=north east,inner sep=0] at ([xshift=17.65cm,yshift=-12.06cm]current page.north west){\color{Classical}\fontsize{7.3}{9.1}\selectfont };
\node[anchor=north east,inner sep=0] at ([xshift=19.05cm,yshift=-12.06cm]current page.north west){\color{Muted}\fontsize{7.7}{9.3}\selectfont 11/32};
\end{tikzpicture}\end{frame}
\begin{lrbox}{\TitleBox}\begin{minipage}[t]{18.2cm}\raggedright\sloppy\color{Ink}\fontsize{20.7}{24}\selectfont \bfseries\color{Navy}Buffer e sensibilità alla scala di osservazione\end{minipage}\end{lrbox}
\begin{lrbox}{\BodyBox}\begin{minipage}[t]{10.25cm}\raggedright\sloppy\color{Ink}\fontsize{12.4}{16.4}\selectfont Un buffer seleziona i siti entro una distanza dichiarata da una posizione di interesse. Cambiando il raggio cambiano sia la numerosità sia la popolazione locale descritta dalla media. La scala dovrebbe derivare dal processo geologico, dalla risoluzione delle misure e dal quesito, evitando che venga scelta soltanto per ottenere un risultato favorevole. Nel caso sintetico si confrontano raggi di 300, 600 e 900 metri attorno al punto (750,750); la numerosità cresce, mentre la media può variare senza una direzione obbligata.\par\vspace{.28cm}{\centering\fontsize{13}{17}\selectfont\(\displaystyle\begin{gathered}B_r(s_0)=\{s:\|s-s_0\|\le r\}\\ \overline z_r=\frac{\sum_i z_i\mathbf1\{s_i\in B_r(s_0)\}}{\sum_i\mathbf1\{s_i\in B_r(s_0)\}}\end{gathered}\)\par}\end{minipage}\end{lrbox}
\begin{lrbox}{\SideBox}\begin{minipage}[t]{7.45cm}\raggedright\sloppy\color{Ink}\fontsize{12.2}{16.1}\selectfont \input{tikz/M11/D04.tex}\par\vspace{.16cm}{\fontsize{9.5}{12.6}\selectfont $s_0=(750,750)$ m; r: raggio; $z_i$: misura nel sito $s_i$. L’indicatore vale 1 entro il buffer, 0 fuori; occorre almeno un sito incluso.}\end{minipage}\end{lrbox}
\begin{lrbox}{\RecallBox}\begin{minipage}[t]{16.45cm}\raggedright\sloppy\color{Ink}\fontsize{10.1}{13.2}\selectfont Una variazione del raggio modifica la domanda territoriale oltre alla precisione della stima.\end{minipage}\end{lrbox}
\MeasureBox{12-TitleBox}{\TitleBox}
\MeasureBox{12-BodyBox}{\BodyBox}
\MeasureBox{12-SideBox}{\SideBox}
\MeasureBox{12-RecallBox}{\RecallBox}
\begin{frame}[plain]\hypertarget{slide12}{}
\bookmark[level=1,page=12]{Buffer e sensibilità alla scala di osservazione}
\begin{tikzpicture}[remember picture,overlay]\node[anchor=north east,inner sep=0] at ([xshift=19.05cm,yshift=-0.16cm]current page.north west){\color{Classical}\fontsize{8}{10}\selectfont\hyperlink{slide2}{Indice}};
\node[anchor=north west,inner sep=0] at ([xshift=0.95cm,yshift=-0.62cm]current page.north west){\usebox{\TitleBox}};
\draw[Classical,line width=1.3pt] ([xshift=.95cm,yshift=-1.68cm]current page.north west)--([xshift=2.25cm,yshift=-1.68cm]current page.north west);\node[anchor=north west,inner sep=0] at ([xshift=0.95cm,yshift=-2.12cm]current page.north west){\usebox{\BodyBox}};
\node[anchor=north west,inner sep=0] at ([xshift=11.6cm,yshift=-2.1cm]current page.north west){\color{Quantum}\bfseries\fontsize{11.5}{14}\selectfont Lettura del diagramma};
\node[anchor=north west,inner sep=0] at ([xshift=11.6cm,yshift=-2.75cm]current page.north west){\usebox{\SideBox}};
\draw[Rule,line width=.5pt] ([xshift=.95cm,yshift=-10.43cm]current page.north west)--([xshift=19.05cm,yshift=-10.43cm]current page.north west);\node[anchor=north west,inner sep=0] at ([xshift=0.95cm,yshift=-10.66cm]current page.north west){\color{Classical}\bfseries\fontsize{9.4}{12}\selectfont Richiami};
\node[anchor=north west,inner sep=0] at ([xshift=2.6cm,yshift=-10.64cm]current page.north west){\usebox{\RecallBox}};
\draw[Rule,line width=.35pt] ([xshift=.95cm,yshift=-11.92cm]current page.north west)--([xshift=19.05cm,yshift=-11.92cm]current page.north west);\node[anchor=north west,inner sep=0] at ([xshift=0.95cm,yshift=-12.06cm]current page.north west){\color{Muted}\fontsize{7.7}{9.3}\selectfont M11 · Statistica spaziale};
\node[anchor=north east,inner sep=0] at ([xshift=17.65cm,yshift=-12.06cm]current page.north west){\color{Classical}\fontsize{7.3}{9.1}\selectfont };
\node[anchor=north east,inner sep=0] at ([xshift=19.05cm,yshift=-12.06cm]current page.north west){\color{Muted}\fontsize{7.7}{9.3}\selectfont 12/32};
\end{tikzpicture}\end{frame}
\begin{lrbox}{\TitleBox}\begin{minipage}[t]{18.2cm}\raggedright\sloppy\color{Ink}\fontsize{20.7}{24}\selectfont \bfseries\color{Navy}Costruzione del vicinato e pareggi geometrici\end{minipage}\end{lrbox}
\begin{lrbox}{\BodyBox}\begin{minipage}[t]{10.25cm}\raggedright\sloppy\color{Ink}\fontsize{12.4}{16.4}\selectfont Un vicinato può comprendere tutti i siti entro una soglia oppure un numero prefissato di siti prossimi. Nel reticolo regolare, più punti possono trovarsi alla stessa distanza minima: scegliere il primo nell’ordine della tabella introduce una convenzione che deve essere dichiarata. Per cercare il vicino più prossimo si esclude anzitutto il punto stesso, sostituendo la diagonale della matrice D con infinito. La soglia di 260 metri include invece tutti i vicini orizzontali e verticali, conservando i pareggi geometrici.\par\vspace{.28cm}{\centering\fontsize{13}{17}\selectfont\(\displaystyle\begin{gathered}a_{ij}=\mathbf1\{0<d_{ij}\le260\}\end{gathered}\)\par}\end{minipage}\end{lrbox}
\begin{lrbox}{\SideBox}\begin{minipage}[t]{7.45cm}\raggedright\sloppy\color{Ink}\fontsize{12.2}{16.1}\selectfont \includegraphics[width=7.45cm]{figure/M11/F03.pdf}\par\vspace{.16cm}{\fontsize{9.5}{12.6}\selectfont I 112 segmenti rappresentano coppie non orientate; i siti hanno due, tre oppure quattro vicini.}\end{minipage}\end{lrbox}
\begin{lrbox}{\RecallBox}\begin{minipage}[t]{16.45cm}\raggedright\sloppy\color{Ink}\fontsize{10.1}{13.2}\selectfont Soglia, metrica, pareggi e autoassociazioni sono parti della definizione del vicinato.\end{minipage}\end{lrbox}
\MeasureBox{13-TitleBox}{\TitleBox}
\MeasureBox{13-BodyBox}{\BodyBox}
\MeasureBox{13-SideBox}{\SideBox}
\MeasureBox{13-RecallBox}{\RecallBox}
\begin{frame}[plain]\hypertarget{slide13}{}
\bookmark[level=0,page=13]{Vicinati e matrici dei pesi}
\bookmark[level=1,page=13]{Costruzione del vicinato e pareggi geometrici}
\begin{tikzpicture}[remember picture,overlay]\node[anchor=north east,inner sep=0] at ([xshift=19.05cm,yshift=-0.16cm]current page.north west){\color{Classical}\fontsize{8}{10}\selectfont\hyperlink{slide2}{Indice}};
\node[anchor=north west,inner sep=0] at ([xshift=0.95cm,yshift=-0.62cm]current page.north west){\usebox{\TitleBox}};
\draw[Classical,line width=1.3pt] ([xshift=.95cm,yshift=-1.68cm]current page.north west)--([xshift=2.25cm,yshift=-1.68cm]current page.north west);\node[anchor=north west,inner sep=0] at ([xshift=0.95cm,yshift=-2.12cm]current page.north west){\usebox{\BodyBox}};
\node[anchor=north west,inner sep=0] at ([xshift=11.6cm,yshift=-2.1cm]current page.north west){\color{Quantum}\bfseries\fontsize{11.5}{14}\selectfont Evidenza grafica};
\node[anchor=north west,inner sep=0] at ([xshift=11.6cm,yshift=-2.75cm]current page.north west){\usebox{\SideBox}};
\draw[Rule,line width=.5pt] ([xshift=.95cm,yshift=-10.43cm]current page.north west)--([xshift=19.05cm,yshift=-10.43cm]current page.north west);\node[anchor=north west,inner sep=0] at ([xshift=0.95cm,yshift=-10.66cm]current page.north west){\color{Classical}\bfseries\fontsize{9.4}{12}\selectfont Richiami};
\node[anchor=north west,inner sep=0] at ([xshift=2.6cm,yshift=-10.64cm]current page.north west){\usebox{\RecallBox}};
\draw[Rule,line width=.35pt] ([xshift=.95cm,yshift=-11.92cm]current page.north west)--([xshift=19.05cm,yshift=-11.92cm]current page.north west);\node[anchor=north west,inner sep=0] at ([xshift=0.95cm,yshift=-12.06cm]current page.north west){\color{Muted}\fontsize{7.7}{9.3}\selectfont M11 · Statistica spaziale};
\node[anchor=north east,inner sep=0] at ([xshift=17.65cm,yshift=-12.06cm]current page.north west){\color{Classical}\fontsize{7.3}{9.1}\selectfont };
\node[anchor=north east,inner sep=0] at ([xshift=19.05cm,yshift=-12.06cm]current page.north west){\color{Muted}\fontsize{7.7}{9.3}\selectfont 13/32};
\end{tikzpicture}\end{frame}
\begin{lrbox}{\TitleBox}\begin{minipage}[t]{18.2cm}\raggedright\sloppy\color{Ink}\fontsize{20.7}{24}\selectfont \bfseries\color{Navy}Normalizzazione e assenza di vicini\end{minipage}\end{lrbox}
\begin{lrbox}{\BodyBox}\begin{minipage}[t]{10.25cm}\raggedright\sloppy\color{Ink}\fontsize{12.4}{16.4}\selectfont La matrice binaria A assegna peso unitario a ciascuna relazione ammessa. Dividendo ogni riga per il numero dei suoi vicini si ottiene W, le cui righe sommano a uno quando il sito possiede almeno un vicino. La trasformazione attribuisce peso totale uguale ai siti, pur producendo pesi individuali maggiori lungo il bordo. Un sito isolato genera un denominatore nullo e richiede una decisione esplicita: modificare il vicinato, escludere il sito con motivazione o adottare una convenzione coerente con l’analisi.\par\vspace{.28cm}{\centering\fontsize{13}{17}\selectfont\(\displaystyle\begin{gathered}w_{ij}=\frac{a_{ij}}{\sum_k a_{ik}},\qquad \sum_jw_{ij}=1\end{gathered}\)\par}\end{minipage}\end{lrbox}
\begin{lrbox}{\SideBox}\begin{minipage}[t]{7.45cm}\raggedright\sloppy\color{Ink}\fontsize{12.2}{16.1}\selectfont \par\vspace{.25cm}\lstinputlisting[language=R]{estratti/M11/pesi.R}\end{minipage}\end{lrbox}
\begin{lrbox}{\RecallBox}\begin{minipage}[t]{16.45cm}\raggedright\sloppy\color{Ink}\fontsize{10.1}{13.2}\selectfont La normalizzazione per riga può trasformare una matrice simmetrica in una matrice asimmetrica.\end{minipage}\end{lrbox}
\MeasureBox{14-TitleBox}{\TitleBox}
\MeasureBox{14-BodyBox}{\BodyBox}
\MeasureBox{14-SideBox}{\SideBox}
\MeasureBox{14-RecallBox}{\RecallBox}
\begin{frame}[plain]\hypertarget{slide14}{}
\bookmark[level=1,page=14]{Normalizzazione e assenza di vicini}
\begin{tikzpicture}[remember picture,overlay]\node[anchor=north east,inner sep=0] at ([xshift=19.05cm,yshift=-0.16cm]current page.north west){\color{Classical}\fontsize{8}{10}\selectfont\hyperlink{slide2}{Indice}};
\node[anchor=north west,inner sep=0] at ([xshift=0.95cm,yshift=-0.62cm]current page.north west){\usebox{\TitleBox}};
\draw[Classical,line width=1.3pt] ([xshift=.95cm,yshift=-1.68cm]current page.north west)--([xshift=2.25cm,yshift=-1.68cm]current page.north west);\node[anchor=north west,inner sep=0] at ([xshift=0.95cm,yshift=-2.12cm]current page.north west){\usebox{\BodyBox}};
\node[anchor=north west,inner sep=0] at ([xshift=11.6cm,yshift=-2.1cm]current page.north west){\color{Quantum}\bfseries\fontsize{11.5}{14}\selectfont Implementazione in R};
\node[anchor=north west,inner sep=0] at ([xshift=11.6cm,yshift=-2.75cm]current page.north west){\usebox{\SideBox}};
\draw[Rule,line width=.5pt] ([xshift=.95cm,yshift=-10.43cm]current page.north west)--([xshift=19.05cm,yshift=-10.43cm]current page.north west);\node[anchor=north west,inner sep=0] at ([xshift=0.95cm,yshift=-10.66cm]current page.north west){\color{Classical}\bfseries\fontsize{9.4}{12}\selectfont Richiami};
\node[anchor=north west,inner sep=0] at ([xshift=2.6cm,yshift=-10.64cm]current page.north west){\usebox{\RecallBox}};
\draw[Rule,line width=.35pt] ([xshift=.95cm,yshift=-11.92cm]current page.north west)--([xshift=19.05cm,yshift=-11.92cm]current page.north west);\node[anchor=north west,inner sep=0] at ([xshift=0.95cm,yshift=-12.06cm]current page.north west){\color{Muted}\fontsize{7.7}{9.3}\selectfont M11 · Statistica spaziale};
\node[anchor=north east,inner sep=0] at ([xshift=17.65cm,yshift=-12.06cm]current page.north west){\color{Classical}\fontsize{7.3}{9.1}\selectfont };
\node[anchor=north east,inner sep=0] at ([xshift=19.05cm,yshift=-12.06cm]current page.north west){\color{Muted}\fontsize{7.7}{9.3}\selectfont 14/32};
\end{tikzpicture}\end{frame}
\begin{lrbox}{\TitleBox}\begin{minipage}[t]{18.2cm}\raggedright\sloppy\color{Ink}\fontsize{20.7}{24}\selectfont \bfseries\color{Navy}Pesi di bordo e media del vicinato\end{minipage}\end{lrbox}
\begin{lrbox}{\BodyBox}\begin{minipage}[t]{10.25cm}\raggedright\sloppy\color{Ink}\fontsize{12.4}{16.4}\selectfont Un sito d’angolo divide il peso fra due vicini, assegnando a ciascuno un mezzo; un sito interno lo divide fra quattro vicini, assegnando un quarto. Perciò una relazione ammessa in entrambe le direzioni può avere pesi differenti nelle due righe di W. Il prodotto Wz restituisce la media locale dei valori contigui, purché l’ordine del vettore coincida con quello delle righe e delle colonne della matrice. Riordinare la sola tabella delle misure comprometterebbe tale corrispondenza senza necessariamente produrre un errore del software.\par\vspace{.28cm}{\centering\fontsize{13}{17}\selectfont\(\displaystyle\begin{gathered}\widetilde z_i=\sum_jw_{ij}z_j\end{gathered}\)\par}\end{minipage}\end{lrbox}
\begin{lrbox}{\SideBox}\begin{minipage}[t]{7.45cm}\raggedright\sloppy\color{Ink}\fontsize{12.2}{16.1}\selectfont \input{tikz/M11/D05.tex}\par\vspace{.16cm}{\fontsize{9.5}{12.6}\selectfont $z_j$: misura nel vicino j; $w_{ij}$: peso dalla riga i. Angolo: due pesi 1/2; interno: quattro pesi 1/4.}\end{minipage}\end{lrbox}
\begin{lrbox}{\RecallBox}\begin{minipage}[t]{16.45cm}\raggedright\sloppy\color{Ink}\fontsize{10.1}{13.2}\selectfont Gli identificativi devono accompagnare matrici, vettori e ogni riordinamento.\end{minipage}\end{lrbox}
\MeasureBox{15-TitleBox}{\TitleBox}
\MeasureBox{15-BodyBox}{\BodyBox}
\MeasureBox{15-SideBox}{\SideBox}
\MeasureBox{15-RecallBox}{\RecallBox}
\begin{frame}[plain]\hypertarget{slide15}{}
\bookmark[level=1,page=15]{Pesi di bordo e media del vicinato}
\begin{tikzpicture}[remember picture,overlay]\node[anchor=north east,inner sep=0] at ([xshift=19.05cm,yshift=-0.16cm]current page.north west){\color{Classical}\fontsize{8}{10}\selectfont\hyperlink{slide2}{Indice}};
\node[anchor=north west,inner sep=0] at ([xshift=0.95cm,yshift=-0.62cm]current page.north west){\usebox{\TitleBox}};
\draw[Classical,line width=1.3pt] ([xshift=.95cm,yshift=-1.68cm]current page.north west)--([xshift=2.25cm,yshift=-1.68cm]current page.north west);\node[anchor=north west,inner sep=0] at ([xshift=0.95cm,yshift=-2.12cm]current page.north west){\usebox{\BodyBox}};
\node[anchor=north west,inner sep=0] at ([xshift=11.6cm,yshift=-2.1cm]current page.north west){\color{Quantum}\bfseries\fontsize{11.5}{14}\selectfont Lettura del diagramma};
\node[anchor=north west,inner sep=0] at ([xshift=11.6cm,yshift=-2.75cm]current page.north west){\usebox{\SideBox}};
\draw[Rule,line width=.5pt] ([xshift=.95cm,yshift=-10.43cm]current page.north west)--([xshift=19.05cm,yshift=-10.43cm]current page.north west);\node[anchor=north west,inner sep=0] at ([xshift=0.95cm,yshift=-10.66cm]current page.north west){\color{Classical}\bfseries\fontsize{9.4}{12}\selectfont Richiami};
\node[anchor=north west,inner sep=0] at ([xshift=2.6cm,yshift=-10.64cm]current page.north west){\usebox{\RecallBox}};
\draw[Rule,line width=.35pt] ([xshift=.95cm,yshift=-11.92cm]current page.north west)--([xshift=19.05cm,yshift=-11.92cm]current page.north west);\node[anchor=north west,inner sep=0] at ([xshift=0.95cm,yshift=-12.06cm]current page.north west){\color{Muted}\fontsize{7.7}{9.3}\selectfont M11 · Statistica spaziale};
\node[anchor=north east,inner sep=0] at ([xshift=17.65cm,yshift=-12.06cm]current page.north west){\color{Classical}\fontsize{7.3}{9.1}\selectfont };
\node[anchor=north east,inner sep=0] at ([xshift=19.05cm,yshift=-12.06cm]current page.north west){\color{Muted}\fontsize{7.7}{9.3}\selectfont 15/32};
\end{tikzpicture}\end{frame}
\begin{lrbox}{\TitleBox}\begin{minipage}[t]{18.2cm}\raggedright\sloppy\color{Ink}\fontsize{20.7}{24}\selectfont \bfseries\color{Navy}L’indice di Moran come prodotto quadratico\end{minipage}\end{lrbox}
\begin{lrbox}{\BodyBox}\begin{minipage}[t]{10.25cm}\raggedright\sloppy\color{Ink}\fontsize{12.4}{16.4}\selectfont L’indice di Moran confronta il prodotto degli scarti nei siti collegati con la dispersione complessiva. Il fattore n diviso S\textsubscript{0} corregge la somma totale dei pesi, permettendo di riconoscere la convenzione utilizzata. Con pesi binari entro 260 metri, il reticolo produce un indice prossimo a 0,720. Sostituendo A con W cambia la ponderazione dei siti di bordo e può cambiare il valore dell’indice. La positività segnala concordanza rispetto al vicinato scelto, senza stabilire da sola quale processo l’abbia generata.\par\vspace{.28cm}{\centering\fontsize{13}{17}\selectfont\(\displaystyle\begin{gathered}I=\frac{n}{S_0}\frac{\sum_i\sum_j a_{ij}(z_i-\overline z)(z_j-\overline z)}{\sum_i(z_i-\overline z)^2},\\ S_0=\sum_{ij}a_{ij}\end{gathered}\)\par}\end{minipage}\end{lrbox}
\begin{lrbox}{\SideBox}\begin{minipage}[t]{7.45cm}\raggedright\sloppy\color{Ink}\fontsize{12.2}{16.1}\selectfont n = 64 siti; $\bar z$ è la media; $a_{ij}$ vale 1 per siti distinti entro 260 m e 0 altrimenti. Il numeratore è il prodotto quadratico degli scarti con A. $S_0=224$, perché ogni coppia non orientata compare due volte. Pesi diversi definiscono confronti spaziali diversi.\end{minipage}\end{lrbox}
\begin{lrbox}{\RecallBox}\begin{minipage}[t]{16.45cm}\raggedright\sloppy\color{Ink}\fontsize{10.1}{13.2}\selectfont L’indice deve essere accompagnato dalla definizione dei pesi e dal dominio campionato.\end{minipage}\end{lrbox}
\MeasureBox{16-TitleBox}{\TitleBox}
\MeasureBox{16-BodyBox}{\BodyBox}
\MeasureBox{16-SideBox}{\SideBox}
\MeasureBox{16-RecallBox}{\RecallBox}
\begin{frame}[plain]\hypertarget{slide16}{}
\bookmark[level=0,page=16]{Associazione spaziale e riferimento inferenziale}
\bookmark[level=1,page=16]{L’indice di Moran come prodotto quadratico}
\begin{tikzpicture}[remember picture,overlay]\node[anchor=north east,inner sep=0] at ([xshift=19.05cm,yshift=-0.16cm]current page.north west){\color{Classical}\fontsize{8}{10}\selectfont\hyperlink{slide2}{Indice}};
\node[anchor=north west,inner sep=0] at ([xshift=0.95cm,yshift=-0.62cm]current page.north west){\usebox{\TitleBox}};
\draw[Classical,line width=1.3pt] ([xshift=.95cm,yshift=-1.68cm]current page.north west)--([xshift=2.25cm,yshift=-1.68cm]current page.north west);\node[anchor=north west,inner sep=0] at ([xshift=0.95cm,yshift=-2.12cm]current page.north west){\usebox{\BodyBox}};
\node[anchor=north west,inner sep=0] at ([xshift=11.6cm,yshift=-2.1cm]current page.north west){\color{Quantum}\bfseries\fontsize{11.5}{14}\selectfont Sviluppo e interpretazione};
\node[anchor=north west,inner sep=0] at ([xshift=11.6cm,yshift=-2.75cm]current page.north west){\usebox{\SideBox}};
\draw[Rule,line width=.5pt] ([xshift=.95cm,yshift=-10.43cm]current page.north west)--([xshift=19.05cm,yshift=-10.43cm]current page.north west);\node[anchor=north west,inner sep=0] at ([xshift=0.95cm,yshift=-10.66cm]current page.north west){\color{Classical}\bfseries\fontsize{9.4}{12}\selectfont Richiami};
\node[anchor=north west,inner sep=0] at ([xshift=2.6cm,yshift=-10.64cm]current page.north west){\usebox{\RecallBox}};
\draw[Rule,line width=.35pt] ([xshift=.95cm,yshift=-11.92cm]current page.north west)--([xshift=19.05cm,yshift=-11.92cm]current page.north west);\node[anchor=north west,inner sep=0] at ([xshift=0.95cm,yshift=-12.06cm]current page.north west){\color{Muted}\fontsize{7.7}{9.3}\selectfont M11 · Statistica spaziale};
\node[anchor=north east,inner sep=0] at ([xshift=17.65cm,yshift=-12.06cm]current page.north west){\color{Classical}\fontsize{7.3}{9.1}\selectfont };
\node[anchor=north east,inner sep=0] at ([xshift=19.05cm,yshift=-12.06cm]current page.north west){\color{Muted}\fontsize{7.7}{9.3}\selectfont 16/32};
\end{tikzpicture}\end{frame}
\begin{lrbox}{\TitleBox}\begin{minipage}[t]{18.2cm}\raggedright\sloppy\color{Ink}\fontsize{20.7}{24}\selectfont \bfseries\color{Navy}Controllo numerico della formula di Moran\end{minipage}\end{lrbox}
\begin{lrbox}{\BodyBox}\begin{minipage}[t]{10.25cm}\raggedright\sloppy\color{Ink}\fontsize{12.4}{16.4}\selectfont Il codice costruisce esplicitamente gli scarti, i prodotti fra coppie e la somma dei pesi, rendendo verificabile ogni termine della formula. Una seconda implementazione mediante prodotto matriciale deve coincidere entro la tolleranza numerica. Questa equivalenza controlla il calcolo, mentre la pertinenza dei pesi rimane una scelta scientifica. Poiché una tendenza territoriale può generare concordanza fra siti vicini, l’indice calcolato sui valori grezzi combina eventuali gradienti sistematici e dipendenza residua.\end{minipage}\end{lrbox}
\begin{lrbox}{\SideBox}\begin{minipage}[t]{7.45cm}\raggedright\sloppy\color{Ink}\fontsize{12.2}{16.1}\selectfont \par\vspace{.25cm}\lstinputlisting[language=R]{estratti/M11/moran.R}\end{minipage}\end{lrbox}
\begin{lrbox}{\RecallBox}\begin{minipage}[t]{16.45cm}\raggedright\sloppy\color{Ink}\fontsize{10.1}{13.2}\selectfont La correttezza dell’implementazione deve precedere l’interpretazione, che richiede ipotesi ulteriori.\end{minipage}\end{lrbox}
\MeasureBox{17-TitleBox}{\TitleBox}
\MeasureBox{17-BodyBox}{\BodyBox}
\MeasureBox{17-SideBox}{\SideBox}
\MeasureBox{17-RecallBox}{\RecallBox}
\begin{frame}[plain]\hypertarget{slide17}{}
\bookmark[level=1,page=17]{Controllo numerico della formula di Moran}
\begin{tikzpicture}[remember picture,overlay]\node[anchor=north east,inner sep=0] at ([xshift=19.05cm,yshift=-0.16cm]current page.north west){\color{Classical}\fontsize{8}{10}\selectfont\hyperlink{slide2}{Indice}};
\node[anchor=north west,inner sep=0] at ([xshift=0.95cm,yshift=-0.62cm]current page.north west){\usebox{\TitleBox}};
\draw[Classical,line width=1.3pt] ([xshift=.95cm,yshift=-1.68cm]current page.north west)--([xshift=2.25cm,yshift=-1.68cm]current page.north west);\node[anchor=north west,inner sep=0] at ([xshift=0.95cm,yshift=-2.12cm]current page.north west){\usebox{\BodyBox}};
\node[anchor=north west,inner sep=0] at ([xshift=11.6cm,yshift=-2.1cm]current page.north west){\color{Quantum}\bfseries\fontsize{11.5}{14}\selectfont Implementazione in R};
\node[anchor=north west,inner sep=0] at ([xshift=11.6cm,yshift=-2.75cm]current page.north west){\usebox{\SideBox}};
\draw[Rule,line width=.5pt] ([xshift=.95cm,yshift=-10.43cm]current page.north west)--([xshift=19.05cm,yshift=-10.43cm]current page.north west);\node[anchor=north west,inner sep=0] at ([xshift=0.95cm,yshift=-10.66cm]current page.north west){\color{Classical}\bfseries\fontsize{9.4}{12}\selectfont Richiami};
\node[anchor=north west,inner sep=0] at ([xshift=2.6cm,yshift=-10.64cm]current page.north west){\usebox{\RecallBox}};
\draw[Rule,line width=.35pt] ([xshift=.95cm,yshift=-11.92cm]current page.north west)--([xshift=19.05cm,yshift=-11.92cm]current page.north west);\node[anchor=north west,inner sep=0] at ([xshift=0.95cm,yshift=-12.06cm]current page.north west){\color{Muted}\fontsize{7.7}{9.3}\selectfont M11 · Statistica spaziale};
\node[anchor=north east,inner sep=0] at ([xshift=17.65cm,yshift=-12.06cm]current page.north west){\color{Classical}\fontsize{7.3}{9.1}\selectfont };
\node[anchor=north east,inner sep=0] at ([xshift=19.05cm,yshift=-12.06cm]current page.north west){\color{Muted}\fontsize{7.7}{9.3}\selectfont 17/32};
\end{tikzpicture}\end{frame}
\begin{lrbox}{\TitleBox}\begin{minipage}[t]{18.2cm}\raggedright\sloppy\color{Ink}\fontsize{20.7}{24}\selectfont \bfseries\color{Navy}Diagramma degli scarti e media dei vicini\end{minipage}\end{lrbox}
\begin{lrbox}{\BodyBox}\begin{minipage}[t]{10.25cm}\raggedright\sloppy\color{Ink}\fontsize{12.4}{16.4}\selectfont Il diagramma mette in relazione lo scarto di ciascun sito con lo scarto medio del suo vicinato. I quadranti alto–alto e basso–basso raccolgono concordanze positive; gli altri due raccolgono discordanza locale. Con W normalizzata per riga e scarti centrati, la pendenza della regressione con intercetta coincide con la forma di Moran basata su W. Tale pendenza può differire dall’indice con pesi binari, poiché la normalizzazione cambia il contributo dei siti. Il diagramma descrive una relazione aggregata e non costituisce un test locale per ogni punto.\par\vspace{.28cm}{\centering\fontsize{13}{17}\selectfont\(\displaystyle\begin{gathered}\widehat\beta=\frac{v^\top Wv}{v^\top v}=I_W\end{gathered}\)\par}\end{minipage}\end{lrbox}
\begin{lrbox}{\SideBox}\begin{minipage}[t]{7.45cm}\raggedright\sloppy\color{Ink}\fontsize{12.2}{16.1}\selectfont \includegraphics[width=7.45cm]{figure/M11/F04.pdf}\par\vspace{.16cm}{\fontsize{9.5}{12.6}\selectfont $v_i=z_i-\bar z$; W: pesi normalizzati per riga. Ogni punto è un sito; le linee a zero delimitano i quattro quadranti.}\end{minipage}\end{lrbox}
\begin{lrbox}{\RecallBox}\begin{minipage}[t]{16.45cm}\raggedright\sloppy\color{Ink}\fontsize{10.1}{13.2}\selectfont Grafico e indice devono impiegare la stessa convenzione dei pesi quando vengono confrontati.\end{minipage}\end{lrbox}
\MeasureBox{18-TitleBox}{\TitleBox}
\MeasureBox{18-BodyBox}{\BodyBox}
\MeasureBox{18-SideBox}{\SideBox}
\MeasureBox{18-RecallBox}{\RecallBox}
\begin{frame}[plain]\hypertarget{slide18}{}
\bookmark[level=1,page=18]{Diagramma degli scarti e media dei vicini}
\begin{tikzpicture}[remember picture,overlay]\node[anchor=north east,inner sep=0] at ([xshift=19.05cm,yshift=-0.16cm]current page.north west){\color{Classical}\fontsize{8}{10}\selectfont\hyperlink{slide2}{Indice}};
\node[anchor=north west,inner sep=0] at ([xshift=0.95cm,yshift=-0.62cm]current page.north west){\usebox{\TitleBox}};
\draw[Classical,line width=1.3pt] ([xshift=.95cm,yshift=-1.68cm]current page.north west)--([xshift=2.25cm,yshift=-1.68cm]current page.north west);\node[anchor=north west,inner sep=0] at ([xshift=0.95cm,yshift=-2.12cm]current page.north west){\usebox{\BodyBox}};
\node[anchor=north west,inner sep=0] at ([xshift=11.6cm,yshift=-2.1cm]current page.north west){\color{Quantum}\bfseries\fontsize{11.5}{14}\selectfont Evidenza grafica};
\node[anchor=north west,inner sep=0] at ([xshift=11.6cm,yshift=-2.75cm]current page.north west){\usebox{\SideBox}};
\draw[Rule,line width=.5pt] ([xshift=.95cm,yshift=-10.43cm]current page.north west)--([xshift=19.05cm,yshift=-10.43cm]current page.north west);\node[anchor=north west,inner sep=0] at ([xshift=0.95cm,yshift=-10.66cm]current page.north west){\color{Classical}\bfseries\fontsize{9.4}{12}\selectfont Richiami};
\node[anchor=north west,inner sep=0] at ([xshift=2.6cm,yshift=-10.64cm]current page.north west){\usebox{\RecallBox}};
\draw[Rule,line width=.35pt] ([xshift=.95cm,yshift=-11.92cm]current page.north west)--([xshift=19.05cm,yshift=-11.92cm]current page.north west);\node[anchor=north west,inner sep=0] at ([xshift=0.95cm,yshift=-12.06cm]current page.north west){\color{Muted}\fontsize{7.7}{9.3}\selectfont M11 · Statistica spaziale};
\node[anchor=north east,inner sep=0] at ([xshift=17.65cm,yshift=-12.06cm]current page.north west){\color{Classical}\fontsize{7.3}{9.1}\selectfont };
\node[anchor=north east,inner sep=0] at ([xshift=19.05cm,yshift=-12.06cm]current page.north west){\color{Muted}\fontsize{7.7}{9.3}\selectfont 18/32};
\end{tikzpicture}\end{frame}
\begin{lrbox}{\TitleBox}\begin{minipage}[t]{18.2cm}\raggedright\sloppy\color{Ink}\fontsize{20.7}{24}\selectfont \bfseries\color{Navy}Permutazioni condizionate alle posizioni\end{minipage}\end{lrbox}
\begin{lrbox}{\BodyBox}\begin{minipage}[t]{10.25cm}\raggedright\sloppy\color{Ink}\fontsize{12.4}{16.4}\selectfont Il riferimento di permutazione mantiene fissi siti, pesi e insieme dei valori, riassegnando questi ultimi alle posizioni. La sua interpretazione richiede che, sotto l’ipotesi nulla, tali assegnazioni siano scambiabili. Si utilizzano 999 permutazioni con seme documentato, ottenendo una risoluzione di un millesimo per il p-value corretto. Il p-value unilaterale conta indici almeno pari a quello osservato, includendo la correzione di una unità che impedisce di attribuire probabilità nulla a un esito mai incontrato.\par\vspace{.28cm}{\centering\fontsize{13}{17}\selectfont\(\displaystyle\begin{gathered}p_{\rm MC}=\frac{1+\sum_{b=1}^B\mathbf1\{I_b\ge I_{\rm obs}\}}{B+1}\end{gathered}\)\par}\end{minipage}\end{lrbox}
\begin{lrbox}{\SideBox}\begin{minipage}[t]{7.45cm}\raggedright\sloppy\color{Ink}\fontsize{12.2}{16.1}\selectfont \includegraphics[width=7.45cm]{figure/M11/F05.pdf}\par\vspace{.16cm}{\fontsize{9.5}{12.6}\selectfont $I_b$: indice alla permutazione b; $I_{\rm obs}$: indice osservato (linea viola). B = 999 permutazioni, con siti e pesi fissi.}\end{minipage}\end{lrbox}
\begin{lrbox}{\RecallBox}\begin{minipage}[t]{16.45cm}\raggedright\sloppy\color{Ink}\fontsize{10.1}{13.2}\selectfont Scambiabilità, direzione del confronto e numero di permutazioni definiscono il riferimento del test.\end{minipage}\end{lrbox}
\MeasureBox{19-TitleBox}{\TitleBox}
\MeasureBox{19-BodyBox}{\BodyBox}
\MeasureBox{19-SideBox}{\SideBox}
\MeasureBox{19-RecallBox}{\RecallBox}
\begin{frame}[plain]\hypertarget{slide19}{}
\bookmark[level=1,page=19]{Permutazioni condizionate alle posizioni}
\begin{tikzpicture}[remember picture,overlay]\node[anchor=north east,inner sep=0] at ([xshift=19.05cm,yshift=-0.16cm]current page.north west){\color{Classical}\fontsize{8}{10}\selectfont\hyperlink{slide2}{Indice}};
\node[anchor=north west,inner sep=0] at ([xshift=0.95cm,yshift=-0.62cm]current page.north west){\usebox{\TitleBox}};
\draw[Classical,line width=1.3pt] ([xshift=.95cm,yshift=-1.68cm]current page.north west)--([xshift=2.25cm,yshift=-1.68cm]current page.north west);\node[anchor=north west,inner sep=0] at ([xshift=0.95cm,yshift=-2.12cm]current page.north west){\usebox{\BodyBox}};
\node[anchor=north west,inner sep=0] at ([xshift=11.6cm,yshift=-2.1cm]current page.north west){\color{Quantum}\bfseries\fontsize{11.5}{14}\selectfont Evidenza grafica};
\node[anchor=north west,inner sep=0] at ([xshift=11.6cm,yshift=-2.75cm]current page.north west){\usebox{\SideBox}};
\draw[Rule,line width=.5pt] ([xshift=.95cm,yshift=-10.43cm]current page.north west)--([xshift=19.05cm,yshift=-10.43cm]current page.north west);\node[anchor=north west,inner sep=0] at ([xshift=0.95cm,yshift=-10.66cm]current page.north west){\color{Classical}\bfseries\fontsize{9.4}{12}\selectfont Richiami};
\node[anchor=north west,inner sep=0] at ([xshift=2.6cm,yshift=-10.64cm]current page.north west){\usebox{\RecallBox}};
\draw[Rule,line width=.35pt] ([xshift=.95cm,yshift=-11.92cm]current page.north west)--([xshift=19.05cm,yshift=-11.92cm]current page.north west);\node[anchor=north west,inner sep=0] at ([xshift=0.95cm,yshift=-12.06cm]current page.north west){\color{Muted}\fontsize{7.7}{9.3}\selectfont M11 · Statistica spaziale};
\node[anchor=north east,inner sep=0] at ([xshift=17.65cm,yshift=-12.06cm]current page.north west){\color{Classical}\fontsize{7.3}{9.1}\selectfont };
\node[anchor=north east,inner sep=0] at ([xshift=19.05cm,yshift=-12.06cm]current page.north west){\color{Muted}\fontsize{7.7}{9.3}\selectfont 19/32};
\end{tikzpicture}\end{frame}
\begin{lrbox}{\TitleBox}\begin{minipage}[t]{18.2cm}\raggedright\sloppy\color{Ink}\fontsize{20.7}{24}\selectfont \bfseries\color{Navy}Tendenza, residui e interpretazione geologica\end{minipage}\end{lrbox}
\begin{lrbox}{\BodyBox}\begin{minipage}[t]{10.25cm}\raggedright\sloppy\color{Ink}\fontsize{12.4}{16.4}\selectfont Una concentrazione che aumenta progressivamente verso Est può produrre un indice di Moran positivo anche quando gli scarti attorno al gradiente sono indipendenti. Se l’obiettivo riguarda l’associazione residua, occorre specificare prima il modello della media e valutare la dipendenza delle sue innovazioni o dei suoi residui. La stima del gradiente modifica la distribuzione del successivo indice, per cui permutare ingenuamente i residui richiede ulteriori giustificazioni. Il quesito geologico stabilisce se interessino il campo complessivo, il gradiente o la struttura rimasta dopo la sua rimozione.\par\vspace{.28cm}{\centering\fontsize{13}{17}\selectfont\(\displaystyle\begin{gathered}Z(s)=m(s)+\varepsilon(s)\end{gathered}\)\par}\end{minipage}\end{lrbox}
\begin{lrbox}{\SideBox}\begin{minipage}[t]{7.45cm}\raggedright\sloppy\color{Ink}\fontsize{12.2}{16.1}\selectfont \input{tikz/M11/D06.tex}\end{minipage}\end{lrbox}
\begin{lrbox}{\RecallBox}\begin{minipage}[t]{16.45cm}\raggedright\sloppy\color{Ink}\fontsize{10.1}{13.2}\selectfont Un indice di associazione spaziale non identifica automaticamente un meccanismo causale.\end{minipage}\end{lrbox}
\MeasureBox{20-TitleBox}{\TitleBox}
\MeasureBox{20-BodyBox}{\BodyBox}
\MeasureBox{20-SideBox}{\SideBox}
\MeasureBox{20-RecallBox}{\RecallBox}
\begin{frame}[plain]\hypertarget{slide20}{}
\bookmark[level=1,page=20]{Tendenza, residui e interpretazione geologica}
\begin{tikzpicture}[remember picture,overlay]\node[anchor=north east,inner sep=0] at ([xshift=19.05cm,yshift=-0.16cm]current page.north west){\color{Classical}\fontsize{8}{10}\selectfont\hyperlink{slide2}{Indice}};
\node[anchor=north west,inner sep=0] at ([xshift=0.95cm,yshift=-0.62cm]current page.north west){\usebox{\TitleBox}};
\draw[Classical,line width=1.3pt] ([xshift=.95cm,yshift=-1.68cm]current page.north west)--([xshift=2.25cm,yshift=-1.68cm]current page.north west);\node[anchor=north west,inner sep=0] at ([xshift=0.95cm,yshift=-2.12cm]current page.north west){\usebox{\BodyBox}};
\node[anchor=north west,inner sep=0] at ([xshift=11.6cm,yshift=-2.1cm]current page.north west){\color{Quantum}\bfseries\fontsize{11.5}{14}\selectfont Lettura del diagramma};
\node[anchor=north west,inner sep=0] at ([xshift=11.6cm,yshift=-2.75cm]current page.north west){\usebox{\SideBox}};
\draw[Rule,line width=.5pt] ([xshift=.95cm,yshift=-10.43cm]current page.north west)--([xshift=19.05cm,yshift=-10.43cm]current page.north west);\node[anchor=north west,inner sep=0] at ([xshift=0.95cm,yshift=-10.66cm]current page.north west){\color{Classical}\bfseries\fontsize{9.4}{12}\selectfont Richiami};
\node[anchor=north west,inner sep=0] at ([xshift=2.6cm,yshift=-10.64cm]current page.north west){\usebox{\RecallBox}};
\draw[Rule,line width=.35pt] ([xshift=.95cm,yshift=-11.92cm]current page.north west)--([xshift=19.05cm,yshift=-11.92cm]current page.north west);\node[anchor=north west,inner sep=0] at ([xshift=0.95cm,yshift=-12.06cm]current page.north west){\color{Muted}\fontsize{7.7}{9.3}\selectfont M11 · Statistica spaziale};
\node[anchor=north east,inner sep=0] at ([xshift=17.65cm,yshift=-12.06cm]current page.north west){\color{Classical}\fontsize{7.3}{9.1}\selectfont };
\node[anchor=north east,inner sep=0] at ([xshift=19.05cm,yshift=-12.06cm]current page.north west){\color{Muted}\fontsize{7.7}{9.3}\selectfont 20/32};
\end{tikzpicture}\end{frame}
\begin{lrbox}{\TitleBox}\begin{minipage}[t]{18.2cm}\raggedright\sloppy\color{Ink}\fontsize{20.7}{24}\selectfont \bfseries\color{Navy}Differenze quadratiche fra coppie di siti\end{minipage}\end{lrbox}
\begin{lrbox}{\BodyBox}\begin{minipage}[t]{10.25cm}\raggedright\sloppy\color{Ink}\fontsize{12.4}{16.4}\selectfont La semivarianza empirica raggruppa coppie di osservazioni secondo la loro distanza e calcola metà della differenza quadratica media. Nel caso sintetico il valore cresce alle distanze maggiori, coerentemente con la presenza di variazioni territoriali ampie. Ogni sito partecipa a molte coppie, cosicché le 2016 differenze non costituiscono altrettante repliche indipendenti. Classi di distanza, numerosità per classe, orientamento e presenza di tendenze devono accompagnare il grafico, del quale condizionano l’interpretazione.\par\vspace{.28cm}{\centering\fontsize{13}{17}\selectfont\(\displaystyle\begin{gathered}\widehat\gamma(h)=\frac{1}{2|N(h)|}\sum_{(i,j)\in N(h)}[z_i-z_j]^2\end{gathered}\)\par}\end{minipage}\end{lrbox}
\begin{lrbox}{\SideBox}\begin{minipage}[t]{7.45cm}\raggedright\sloppy\color{Ink}\fontsize{12.2}{16.1}\selectfont \includegraphics[width=7.45cm]{figure/M11/F06.pdf}\par\vspace{.16cm}{\fontsize{9.5}{12.6}\selectfont N(h): coppie nella classe di distanza h; $|N(h)|$: loro numero. Quattro classi, senza duplicati; numerosità esportate in R.}\end{minipage}\end{lrbox}
\begin{lrbox}{\RecallBox}\begin{minipage}[t]{16.45cm}\raggedright\sloppy\color{Ink}\fontsize{10.1}{13.2}\selectfont Le coppie condividono osservazioni e non moltiplicano il numero di unità indipendenti.\end{minipage}\end{lrbox}
\MeasureBox{21-TitleBox}{\TitleBox}
\MeasureBox{21-BodyBox}{\BodyBox}
\MeasureBox{21-SideBox}{\SideBox}
\MeasureBox{21-RecallBox}{\RecallBox}
\begin{frame}[plain]\hypertarget{slide21}{}
\bookmark[level=0,page=21]{Semivarianza e struttura a distanza}
\bookmark[level=1,page=21]{Differenze quadratiche fra coppie di siti}
\begin{tikzpicture}[remember picture,overlay]\node[anchor=north east,inner sep=0] at ([xshift=19.05cm,yshift=-0.16cm]current page.north west){\color{Classical}\fontsize{8}{10}\selectfont\hyperlink{slide2}{Indice}};
\node[anchor=north west,inner sep=0] at ([xshift=0.95cm,yshift=-0.62cm]current page.north west){\usebox{\TitleBox}};
\draw[Classical,line width=1.3pt] ([xshift=.95cm,yshift=-1.68cm]current page.north west)--([xshift=2.25cm,yshift=-1.68cm]current page.north west);\node[anchor=north west,inner sep=0] at ([xshift=0.95cm,yshift=-2.12cm]current page.north west){\usebox{\BodyBox}};
\node[anchor=north west,inner sep=0] at ([xshift=11.6cm,yshift=-2.1cm]current page.north west){\color{Quantum}\bfseries\fontsize{11.5}{14}\selectfont Evidenza grafica};
\node[anchor=north west,inner sep=0] at ([xshift=11.6cm,yshift=-2.75cm]current page.north west){\usebox{\SideBox}};
\draw[Rule,line width=.5pt] ([xshift=.95cm,yshift=-10.43cm]current page.north west)--([xshift=19.05cm,yshift=-10.43cm]current page.north west);\node[anchor=north west,inner sep=0] at ([xshift=0.95cm,yshift=-10.66cm]current page.north west){\color{Classical}\bfseries\fontsize{9.4}{12}\selectfont Richiami};
\node[anchor=north west,inner sep=0] at ([xshift=2.6cm,yshift=-10.64cm]current page.north west){\usebox{\RecallBox}};
\draw[Rule,line width=.35pt] ([xshift=.95cm,yshift=-11.92cm]current page.north west)--([xshift=19.05cm,yshift=-11.92cm]current page.north west);\node[anchor=north west,inner sep=0] at ([xshift=0.95cm,yshift=-12.06cm]current page.north west){\color{Muted}\fontsize{7.7}{9.3}\selectfont M11 · Statistica spaziale};
\node[anchor=north east,inner sep=0] at ([xshift=17.65cm,yshift=-12.06cm]current page.north west){\color{Classical}\fontsize{7.3}{9.1}\selectfont };
\node[anchor=north east,inner sep=0] at ([xshift=19.05cm,yshift=-12.06cm]current page.north west){\color{Muted}\fontsize{7.7}{9.3}\selectfont 21/32};
\end{tikzpicture}\end{frame}
\begin{lrbox}{\TitleBox}\begin{minipage}[t]{18.2cm}\raggedright\sloppy\color{Ink}\fontsize{20.7}{24}\selectfont \bfseries\color{Navy}Covarianza e variogramma sotto stazionarietà\end{minipage}\end{lrbox}
\begin{lrbox}{\BodyBox}\begin{minipage}[t]{10.25cm}\raggedright\sloppy\color{Ink}\fontsize{12.4}{16.4}\selectfont Per un campo di secondo ordine stazionario, la covarianza dipende dallo spostamento fra due siti e la semivarianza è la differenza fra varianza marginale e covarianza. L’isotropia aggiunge l’ipotesi che conti soltanto la lunghezza dello spostamento, ignorandone la direzione. Un gradiente della media può violare la prima condizione; strutture geologiche orientate possono rendere inadeguata la seconda. Il variogramma descrittivo del reticolo viene pertanto presentato prima di qualunque adattamento di un modello stazionario.\par\vspace{.28cm}{\centering\fontsize{13}{17}\selectfont\(\displaystyle\begin{gathered}\gamma(h)=\tfrac12\operatorname{Var}[Z(s+h)-Z(s)]\ =C(0)-C(h)\end{gathered}\)\par}\end{minipage}\end{lrbox}
\begin{lrbox}{\SideBox}\begin{minipage}[t]{7.45cm}\raggedright\sloppy\color{Ink}\fontsize{12.2}{16.1}\selectfont \input{tikz/M11/D07.tex}\par\vspace{.16cm}{\fontsize{9.5}{12.6}\selectfont h: spostamento vettoriale; C(h): covarianza; C(0): varianza marginale. L’isotropia riduce la dipendenza alla lunghezza di h.}\end{minipage}\end{lrbox}
\begin{lrbox}{\RecallBox}\begin{minipage}[t]{16.45cm}\raggedright\sloppy\color{Ink}\fontsize{10.1}{13.2}\selectfont Stazionarietà e isotropia sono ipotesi del modello, da confrontare con la struttura geologica.\end{minipage}\end{lrbox}
\MeasureBox{22-TitleBox}{\TitleBox}
\MeasureBox{22-BodyBox}{\BodyBox}
\MeasureBox{22-SideBox}{\SideBox}
\MeasureBox{22-RecallBox}{\RecallBox}
\begin{frame}[plain]\hypertarget{slide22}{}
\bookmark[level=1,page=22]{Covarianza e variogramma sotto stazionarietà}
\begin{tikzpicture}[remember picture,overlay]\node[anchor=north east,inner sep=0] at ([xshift=19.05cm,yshift=-0.16cm]current page.north west){\color{Classical}\fontsize{8}{10}\selectfont\hyperlink{slide2}{Indice}};
\node[anchor=north west,inner sep=0] at ([xshift=0.95cm,yshift=-0.62cm]current page.north west){\usebox{\TitleBox}};
\draw[Classical,line width=1.3pt] ([xshift=.95cm,yshift=-1.68cm]current page.north west)--([xshift=2.25cm,yshift=-1.68cm]current page.north west);\node[anchor=north west,inner sep=0] at ([xshift=0.95cm,yshift=-2.12cm]current page.north west){\usebox{\BodyBox}};
\node[anchor=north west,inner sep=0] at ([xshift=11.6cm,yshift=-2.1cm]current page.north west){\color{Quantum}\bfseries\fontsize{11.5}{14}\selectfont Lettura del diagramma};
\node[anchor=north west,inner sep=0] at ([xshift=11.6cm,yshift=-2.75cm]current page.north west){\usebox{\SideBox}};
\draw[Rule,line width=.5pt] ([xshift=.95cm,yshift=-10.43cm]current page.north west)--([xshift=19.05cm,yshift=-10.43cm]current page.north west);\node[anchor=north west,inner sep=0] at ([xshift=0.95cm,yshift=-10.66cm]current page.north west){\color{Classical}\bfseries\fontsize{9.4}{12}\selectfont Richiami};
\node[anchor=north west,inner sep=0] at ([xshift=2.6cm,yshift=-10.64cm]current page.north west){\usebox{\RecallBox}};
\draw[Rule,line width=.35pt] ([xshift=.95cm,yshift=-11.92cm]current page.north west)--([xshift=19.05cm,yshift=-11.92cm]current page.north west);\node[anchor=north west,inner sep=0] at ([xshift=0.95cm,yshift=-12.06cm]current page.north west){\color{Muted}\fontsize{7.7}{9.3}\selectfont M11 · Statistica spaziale};
\node[anchor=north east,inner sep=0] at ([xshift=17.65cm,yshift=-12.06cm]current page.north west){\color{Classical}\fontsize{7.3}{9.1}\selectfont Spatial Data Science, cap. 12};
\node[anchor=north east,inner sep=0] at ([xshift=19.05cm,yshift=-12.06cm]current page.north west){\color{Muted}\fontsize{7.7}{9.3}\selectfont 22/32};
\end{tikzpicture}\end{frame}
\begin{lrbox}{\TitleBox}\begin{minipage}[t]{18.2cm}\raggedright\sloppy\color{Ink}\fontsize{20.7}{24}\selectfont \bfseries\color{Navy}Campionamento di convenienza e copertura\end{minipage}\end{lrbox}
\begin{lrbox}{\BodyBox}\begin{minipage}[t]{10.25cm}\raggedright\sloppy\color{Ink}\fontsize{12.4}{16.4}\selectfont La selezione dei sedici siti lungo la strada fittizia, definita da y minore o uguale a 250 metri, conserva una buona copertura longitudinale e trascura gran parte della direzione Nord–Sud. Repliche precise entro questa fascia ridurrebbero l’incertezza locale senza ristabilire la rappresentatività dell’intero territorio. La carta del disegno permette di discutere quali ambienti siano assenti e quali ipotesi consentirebbero il trasferimento della sintesi. La differenza fra media della fascia e media del reticolo è un esito del caso sintetico, non una stima generale della distorsione di ogni campionamento stradale.\end{minipage}\end{lrbox}
\begin{lrbox}{\SideBox}\begin{minipage}[t]{7.45cm}\raggedright\sloppy\color{Ink}\fontsize{12.2}{16.1}\selectfont \includegraphics[width=7.45cm]{figure/M11/F07.pdf}\par\vspace{.16cm}{\fontsize{9.5}{12.6}\selectfont Viola: sedici siti prossimi alla strada fittizia; grigio: quarantotto siti esclusi dalla selezione.}\end{minipage}\end{lrbox}
\begin{lrbox}{\RecallBox}\begin{minipage}[t]{16.45cm}\raggedright\sloppy\color{Ink}\fontsize{10.1}{13.2}\selectfont La precisione entro il dominio osservato e la rappresentatività territoriale riguardano proprietà differenti.\end{minipage}\end{lrbox}
\MeasureBox{23-TitleBox}{\TitleBox}
\MeasureBox{23-BodyBox}{\BodyBox}
\MeasureBox{23-SideBox}{\SideBox}
\MeasureBox{23-RecallBox}{\RecallBox}
\begin{frame}[plain]\hypertarget{slide23}{}
\bookmark[level=0,page=23]{Copertura e unità della sintesi}
\bookmark[level=1,page=23]{Campionamento di convenienza e copertura}
\begin{tikzpicture}[remember picture,overlay]\node[anchor=north east,inner sep=0] at ([xshift=19.05cm,yshift=-0.16cm]current page.north west){\color{Classical}\fontsize{8}{10}\selectfont\hyperlink{slide2}{Indice}};
\node[anchor=north west,inner sep=0] at ([xshift=0.95cm,yshift=-0.62cm]current page.north west){\usebox{\TitleBox}};
\draw[Classical,line width=1.3pt] ([xshift=.95cm,yshift=-1.68cm]current page.north west)--([xshift=2.25cm,yshift=-1.68cm]current page.north west);\node[anchor=north west,inner sep=0] at ([xshift=0.95cm,yshift=-2.12cm]current page.north west){\usebox{\BodyBox}};
\node[anchor=north west,inner sep=0] at ([xshift=11.6cm,yshift=-2.1cm]current page.north west){\color{Quantum}\bfseries\fontsize{11.5}{14}\selectfont Evidenza grafica};
\node[anchor=north west,inner sep=0] at ([xshift=11.6cm,yshift=-2.75cm]current page.north west){\usebox{\SideBox}};
\draw[Rule,line width=.5pt] ([xshift=.95cm,yshift=-10.43cm]current page.north west)--([xshift=19.05cm,yshift=-10.43cm]current page.north west);\node[anchor=north west,inner sep=0] at ([xshift=0.95cm,yshift=-10.66cm]current page.north west){\color{Classical}\bfseries\fontsize{9.4}{12}\selectfont Richiami};
\node[anchor=north west,inner sep=0] at ([xshift=2.6cm,yshift=-10.64cm]current page.north west){\usebox{\RecallBox}};
\draw[Rule,line width=.35pt] ([xshift=.95cm,yshift=-11.92cm]current page.north west)--([xshift=19.05cm,yshift=-11.92cm]current page.north west);\node[anchor=north west,inner sep=0] at ([xshift=0.95cm,yshift=-12.06cm]current page.north west){\color{Muted}\fontsize{7.7}{9.3}\selectfont M11 · Statistica spaziale};
\node[anchor=north east,inner sep=0] at ([xshift=17.65cm,yshift=-12.06cm]current page.north west){\color{Classical}\fontsize{7.3}{9.1}\selectfont };
\node[anchor=north east,inner sep=0] at ([xshift=19.05cm,yshift=-12.06cm]current page.north west){\color{Muted}\fontsize{7.7}{9.3}\selectfont 23/32};
\end{tikzpicture}\end{frame}
\begin{lrbox}{\TitleBox}\begin{minipage}[t]{18.2cm}\raggedright\sloppy\color{Ink}\fontsize{20.7}{24}\selectfont \bfseries\color{Navy}Medie di blocco e scelta dei pesi\end{minipage}\end{lrbox}
\begin{lrbox}{\BodyBox}\begin{minipage}[t]{10.25cm}\raggedright\sloppy\color{Ink}\fontsize{12.4}{16.4}\selectfont Suddividendo il reticolo in quattro blocchi da sedici siti, la media delle medie di blocco coincide con la media complessiva. Tale uguaglianza deriva dal bilanciamento e non è garantita quando le numerosità differiscono. Per due blocchi con medie 10 e 20 e numerosità 2 e 8, la media dei campioni vale 18, mentre la media equiponderata dei blocchi vale 15. La prima descrive i campioni, la seconda i blocchi equiponderati; una media areale richiederebbe pesi proporzionali alle superfici.\par\vspace{.28cm}{\centering\fontsize{13}{17}\selectfont\(\displaystyle\begin{gathered}\overline z_{\rm campioni}=\frac{\sum_b n_b\overline z_b}{\sum_b n_b}\\ \overline z_{\rm blocchi}=\frac1K\sum_b\overline z_b\end{gathered}\)\par}\end{minipage}\end{lrbox}
\begin{lrbox}{\SideBox}\begin{minipage}[t]{7.45cm}\raggedright\sloppy\color{Ink}\fontsize{12.2}{16.1}\selectfont \input{tikz/M11/D08.tex}\par\vspace{.16cm}{\fontsize{9.5}{12.6}\selectfont K: numero di blocchi; $n_b$: siti nel blocco b; $\bar z_b$: media del blocco. L’esempio confronta pesi per campione e per blocco.}\end{minipage}\end{lrbox}
\begin{lrbox}{\RecallBox}\begin{minipage}[t]{16.45cm}\raggedright\sloppy\color{Ink}\fontsize{10.1}{13.2}\selectfont L’aggregazione deve conservare numerosità, superfici e definizione della popolazione di riferimento.\end{minipage}\end{lrbox}
\MeasureBox{24-TitleBox}{\TitleBox}
\MeasureBox{24-BodyBox}{\BodyBox}
\MeasureBox{24-SideBox}{\SideBox}
\MeasureBox{24-RecallBox}{\RecallBox}
\begin{frame}[plain]\hypertarget{slide24}{}
\bookmark[level=1,page=24]{Medie di blocco e scelta dei pesi}
\begin{tikzpicture}[remember picture,overlay]\node[anchor=north east,inner sep=0] at ([xshift=19.05cm,yshift=-0.16cm]current page.north west){\color{Classical}\fontsize{8}{10}\selectfont\hyperlink{slide2}{Indice}};
\node[anchor=north west,inner sep=0] at ([xshift=0.95cm,yshift=-0.62cm]current page.north west){\usebox{\TitleBox}};
\draw[Classical,line width=1.3pt] ([xshift=.95cm,yshift=-1.68cm]current page.north west)--([xshift=2.25cm,yshift=-1.68cm]current page.north west);\node[anchor=north west,inner sep=0] at ([xshift=0.95cm,yshift=-2.12cm]current page.north west){\usebox{\BodyBox}};
\node[anchor=north west,inner sep=0] at ([xshift=11.6cm,yshift=-2.1cm]current page.north west){\color{Quantum}\bfseries\fontsize{11.5}{14}\selectfont Lettura del diagramma};
\node[anchor=north west,inner sep=0] at ([xshift=11.6cm,yshift=-2.75cm]current page.north west){\usebox{\SideBox}};
\draw[Rule,line width=.5pt] ([xshift=.95cm,yshift=-10.43cm]current page.north west)--([xshift=19.05cm,yshift=-10.43cm]current page.north west);\node[anchor=north west,inner sep=0] at ([xshift=0.95cm,yshift=-10.66cm]current page.north west){\color{Classical}\bfseries\fontsize{9.4}{12}\selectfont Richiami};
\node[anchor=north west,inner sep=0] at ([xshift=2.6cm,yshift=-10.64cm]current page.north west){\usebox{\RecallBox}};
\draw[Rule,line width=.35pt] ([xshift=.95cm,yshift=-11.92cm]current page.north west)--([xshift=19.05cm,yshift=-11.92cm]current page.north west);\node[anchor=north west,inner sep=0] at ([xshift=0.95cm,yshift=-12.06cm]current page.north west){\color{Muted}\fontsize{7.7}{9.3}\selectfont M11 · Statistica spaziale};
\node[anchor=north east,inner sep=0] at ([xshift=17.65cm,yshift=-12.06cm]current page.north west){\color{Classical}\fontsize{7.3}{9.1}\selectfont };
\node[anchor=north east,inner sep=0] at ([xshift=19.05cm,yshift=-12.06cm]current page.north west){\color{Muted}\fontsize{7.7}{9.3}\selectfont 24/32};
\end{tikzpicture}\end{frame}
\begin{lrbox}{\TitleBox}\begin{minipage}[t]{18.2cm}\raggedright\sloppy\color{Ink}\fontsize{20.7}{24}\selectfont \bfseries\color{Navy}Interpolazione a distanza inversa\end{minipage}\end{lrbox}
\begin{lrbox}{\BodyBox}\begin{minipage}[t]{10.25cm}\raggedright\sloppy\color{Ink}\fontsize{12.4}{16.4}\selectfont L’interpolazione IDW stima il valore in una posizione priva di misura mediante una media pesata delle osservazioni. I pesi diminuiscono come una potenza della distanza e vengono normalizzati affinché sommino a uno. Aumentando l’esponente si concentra il contributo sui siti prossimi, mentre esponenti più piccoli distribuiscono l’informazione su un’area più ampia. La regola deve essere accompagnata da una scelta di metrica, vicinato e trattamento dei siti coincidenti. Il punto di esempio è situato a (900,900) metri, fra i nodi del reticolo.\par\vspace{.28cm}{\centering\fontsize{13}{17}\selectfont\(\displaystyle\begin{gathered}\widehat z(s_0)=\sum_i\lambda_i(s_0)z_i,\qquad \lambda_i=\frac{d_{i0}^{-p}}{\sum_jd_{j0}^{-p}}\end{gathered}\)\par}\end{minipage}\end{lrbox}
\begin{lrbox}{\SideBox}\begin{minipage}[t]{7.45cm}\raggedright\sloppy\color{Ink}\fontsize{12.2}{16.1}\selectfont \includegraphics[width=7.45cm]{figure/M11/F10.pdf}\par\vspace{.16cm}{\fontsize{9.5}{12.6}\selectfont $d_{i0}$: distanza dal sito i a (900,900) m; p > 0: potenza. Curve dei pesi per p = 1, 2, 3; tutte le distanze sono positive.}\end{minipage}\end{lrbox}
\begin{lrbox}{\RecallBox}\begin{minipage}[t]{16.45cm}\raggedright\sloppy\color{Ink}\fontsize{10.1}{13.2}\selectfont La potenza controlla la località della media e deve essere valutata rispetto al compito predittivo.\end{minipage}\end{lrbox}
\MeasureBox{25-TitleBox}{\TitleBox}
\MeasureBox{25-BodyBox}{\BodyBox}
\MeasureBox{25-SideBox}{\SideBox}
\MeasureBox{25-RecallBox}{\RecallBox}
\begin{frame}[plain]\hypertarget{slide25}{}
\bookmark[level=0,page=25]{Interpolazione e verifica geografica}
\bookmark[level=1,page=25]{Interpolazione a distanza inversa}
\begin{tikzpicture}[remember picture,overlay]\node[anchor=north east,inner sep=0] at ([xshift=19.05cm,yshift=-0.16cm]current page.north west){\color{Classical}\fontsize{8}{10}\selectfont\hyperlink{slide2}{Indice}};
\node[anchor=north west,inner sep=0] at ([xshift=0.95cm,yshift=-0.62cm]current page.north west){\usebox{\TitleBox}};
\draw[Classical,line width=1.3pt] ([xshift=.95cm,yshift=-1.68cm]current page.north west)--([xshift=2.25cm,yshift=-1.68cm]current page.north west);\node[anchor=north west,inner sep=0] at ([xshift=0.95cm,yshift=-2.12cm]current page.north west){\usebox{\BodyBox}};
\node[anchor=north west,inner sep=0] at ([xshift=11.6cm,yshift=-2.1cm]current page.north west){\color{Quantum}\bfseries\fontsize{11.5}{14}\selectfont Evidenza grafica};
\node[anchor=north west,inner sep=0] at ([xshift=11.6cm,yshift=-2.75cm]current page.north west){\usebox{\SideBox}};
\draw[Rule,line width=.5pt] ([xshift=.95cm,yshift=-10.43cm]current page.north west)--([xshift=19.05cm,yshift=-10.43cm]current page.north west);\node[anchor=north west,inner sep=0] at ([xshift=0.95cm,yshift=-10.66cm]current page.north west){\color{Classical}\bfseries\fontsize{9.4}{12}\selectfont Richiami};
\node[anchor=north west,inner sep=0] at ([xshift=2.6cm,yshift=-10.64cm]current page.north west){\usebox{\RecallBox}};
\draw[Rule,line width=.35pt] ([xshift=.95cm,yshift=-11.92cm]current page.north west)--([xshift=19.05cm,yshift=-11.92cm]current page.north west);\node[anchor=north west,inner sep=0] at ([xshift=0.95cm,yshift=-12.06cm]current page.north west){\color{Muted}\fontsize{7.7}{9.3}\selectfont M11 · Statistica spaziale};
\node[anchor=north east,inner sep=0] at ([xshift=17.65cm,yshift=-12.06cm]current page.north west){\color{Classical}\fontsize{7.3}{9.1}\selectfont };
\node[anchor=north east,inner sep=0] at ([xshift=19.05cm,yshift=-12.06cm]current page.north west){\color{Muted}\fontsize{7.7}{9.3}\selectfont 25/32};
\end{tikzpicture}\end{frame}
\begin{lrbox}{\TitleBox}\begin{minipage}[t]{18.2cm}\raggedright\sloppy\color{Ink}\fontsize{20.7}{24}\selectfont \bfseries\color{Navy}Vincoli della media pesata e siti coincidenti\end{minipage}\end{lrbox}
\begin{lrbox}{\BodyBox}\begin{minipage}[t]{10.25cm}\raggedright\sloppy\color{Ink}\fontsize{12.4}{16.4}\selectfont La positività dei pesi normalizzati mantiene ogni previsione IDW entro l’intervallo dei valori utilizzati. Questa proprietà offre un controllo numerico e delimita la capacità del metodo di rappresentare estremi non osservati. Se la posizione richiesta coincide con un sito campionato, la distanza nulla deve essere trattata prima dell’elevamento a potenza negativa. Con una sola misura si restituisce il valore del sito; con repliche discordanti occorre definire un modello o una regola di aggregazione, poiché l’interpolazione esatta di entrambe sarebbe impossibile.\par\vspace{.28cm}{\centering\fontsize{13}{17}\selectfont\(\displaystyle\begin{gathered}\min_i z_i\le\widehat z(s_0)\le\max_i z_i\end{gathered}\)\par}\end{minipage}\end{lrbox}
\begin{lrbox}{\SideBox}\begin{minipage}[t]{7.45cm}\raggedright\sloppy\color{Ink}\fontsize{12.2}{16.1}\selectfont \par\vspace{.25cm}\lstinputlisting[language=R]{estratti/M11/idw.R}\end{minipage}\end{lrbox}
\begin{lrbox}{\RecallBox}\begin{minipage}[t]{16.45cm}\raggedright\sloppy\color{Ink}\fontsize{10.1}{13.2}\selectfont La gestione delle coincidenze appartiene al modello dei dati e deve precedere il calcolo dei pesi.\end{minipage}\end{lrbox}
\MeasureBox{26-TitleBox}{\TitleBox}
\MeasureBox{26-BodyBox}{\BodyBox}
\MeasureBox{26-SideBox}{\SideBox}
\MeasureBox{26-RecallBox}{\RecallBox}
\begin{frame}[plain]\hypertarget{slide26}{}
\bookmark[level=1,page=26]{Vincoli della media pesata e siti coincidenti}
\begin{tikzpicture}[remember picture,overlay]\node[anchor=north east,inner sep=0] at ([xshift=19.05cm,yshift=-0.16cm]current page.north west){\color{Classical}\fontsize{8}{10}\selectfont\hyperlink{slide2}{Indice}};
\node[anchor=north west,inner sep=0] at ([xshift=0.95cm,yshift=-0.62cm]current page.north west){\usebox{\TitleBox}};
\draw[Classical,line width=1.3pt] ([xshift=.95cm,yshift=-1.68cm]current page.north west)--([xshift=2.25cm,yshift=-1.68cm]current page.north west);\node[anchor=north west,inner sep=0] at ([xshift=0.95cm,yshift=-2.12cm]current page.north west){\usebox{\BodyBox}};
\node[anchor=north west,inner sep=0] at ([xshift=11.6cm,yshift=-2.1cm]current page.north west){\color{Quantum}\bfseries\fontsize{11.5}{14}\selectfont Implementazione in R};
\node[anchor=north west,inner sep=0] at ([xshift=11.6cm,yshift=-2.75cm]current page.north west){\usebox{\SideBox}};
\draw[Rule,line width=.5pt] ([xshift=.95cm,yshift=-10.43cm]current page.north west)--([xshift=19.05cm,yshift=-10.43cm]current page.north west);\node[anchor=north west,inner sep=0] at ([xshift=0.95cm,yshift=-10.66cm]current page.north west){\color{Classical}\bfseries\fontsize{9.4}{12}\selectfont Richiami};
\node[anchor=north west,inner sep=0] at ([xshift=2.6cm,yshift=-10.64cm]current page.north west){\usebox{\RecallBox}};
\draw[Rule,line width=.35pt] ([xshift=.95cm,yshift=-11.92cm]current page.north west)--([xshift=19.05cm,yshift=-11.92cm]current page.north west);\node[anchor=north west,inner sep=0] at ([xshift=0.95cm,yshift=-12.06cm]current page.north west){\color{Muted}\fontsize{7.7}{9.3}\selectfont M11 · Statistica spaziale};
\node[anchor=north east,inner sep=0] at ([xshift=17.65cm,yshift=-12.06cm]current page.north west){\color{Classical}\fontsize{7.3}{9.1}\selectfont };
\node[anchor=north east,inner sep=0] at ([xshift=19.05cm,yshift=-12.06cm]current page.north west){\color{Muted}\fontsize{7.7}{9.3}\selectfont 26/32};
\end{tikzpicture}\end{frame}
\begin{lrbox}{\TitleBox}\begin{minipage}[t]{18.2cm}\raggedright\sloppy\color{Ink}\fontsize{20.7}{24}\selectfont \bfseries\color{Navy}Superficie interpolata e provenienza dei dati\end{minipage}\end{lrbox}
\begin{lrbox}{\BodyBox}\begin{minipage}[t]{10.25cm}\raggedright\sloppy\color{Ink}\fontsize{12.4}{16.4}\selectfont Applicando IDW con potenza due a un reticolo più fitto si ottiene una superficie continua di valori previsti, sulla quale vengono sovrapposti i siti originari. La maggiore densità delle celle migliora la rappresentazione grafica della regola e lascia invariata l’informazione osservativa. Nei settori poco coperti, una superficie visivamente regolare può nascondere l’ampia distanza dalle misure. La carta deve quindi indicare metodo, parametri, unità, copertura e data dell’analisi, affinché il lettore possa distinguere il dato acquisito dalla previsione.\end{minipage}\end{lrbox}
\begin{lrbox}{\SideBox}\begin{minipage}[t]{7.45cm}\raggedright\sloppy\color{Ink}\fontsize{12.2}{16.1}\selectfont \includegraphics[width=7.45cm]{figure/M11/F08.pdf}\par\vspace{.16cm}{\fontsize{9.5}{12.6}\selectfont IDW di potenza due; siti misurati e griglia di 45 × 45 celle.}\end{minipage}\end{lrbox}
\begin{lrbox}{\RecallBox}\begin{minipage}[t]{16.45cm}\raggedright\sloppy\color{Ink}\fontsize{10.1}{13.2}\selectfont La risoluzione grafica di una superficie non coincide con la risoluzione informativa del campionamento.\end{minipage}\end{lrbox}
\MeasureBox{27-TitleBox}{\TitleBox}
\MeasureBox{27-BodyBox}{\BodyBox}
\MeasureBox{27-SideBox}{\SideBox}
\MeasureBox{27-RecallBox}{\RecallBox}
\begin{frame}[plain]\hypertarget{slide27}{}
\bookmark[level=1,page=27]{Superficie interpolata e provenienza dei dati}
\begin{tikzpicture}[remember picture,overlay]\node[anchor=north east,inner sep=0] at ([xshift=19.05cm,yshift=-0.16cm]current page.north west){\color{Classical}\fontsize{8}{10}\selectfont\hyperlink{slide2}{Indice}};
\node[anchor=north west,inner sep=0] at ([xshift=0.95cm,yshift=-0.62cm]current page.north west){\usebox{\TitleBox}};
\draw[Classical,line width=1.3pt] ([xshift=.95cm,yshift=-1.68cm]current page.north west)--([xshift=2.25cm,yshift=-1.68cm]current page.north west);\node[anchor=north west,inner sep=0] at ([xshift=0.95cm,yshift=-2.12cm]current page.north west){\usebox{\BodyBox}};
\node[anchor=north west,inner sep=0] at ([xshift=11.6cm,yshift=-2.1cm]current page.north west){\color{Quantum}\bfseries\fontsize{11.5}{14}\selectfont Evidenza grafica};
\node[anchor=north west,inner sep=0] at ([xshift=11.6cm,yshift=-2.75cm]current page.north west){\usebox{\SideBox}};
\draw[Rule,line width=.5pt] ([xshift=.95cm,yshift=-10.43cm]current page.north west)--([xshift=19.05cm,yshift=-10.43cm]current page.north west);\node[anchor=north west,inner sep=0] at ([xshift=0.95cm,yshift=-10.66cm]current page.north west){\color{Classical}\bfseries\fontsize{9.4}{12}\selectfont Richiami};
\node[anchor=north west,inner sep=0] at ([xshift=2.6cm,yshift=-10.64cm]current page.north west){\usebox{\RecallBox}};
\draw[Rule,line width=.35pt] ([xshift=.95cm,yshift=-11.92cm]current page.north west)--([xshift=19.05cm,yshift=-11.92cm]current page.north west);\node[anchor=north west,inner sep=0] at ([xshift=0.95cm,yshift=-12.06cm]current page.north west){\color{Muted}\fontsize{7.7}{9.3}\selectfont M11 · Statistica spaziale};
\node[anchor=north east,inner sep=0] at ([xshift=17.65cm,yshift=-12.06cm]current page.north west){\color{Classical}\fontsize{7.3}{9.1}\selectfont };
\node[anchor=north east,inner sep=0] at ([xshift=19.05cm,yshift=-12.06cm]current page.north west){\color{Muted}\fontsize{7.7}{9.3}\selectfont 27/32};
\end{tikzpicture}\end{frame}
\begin{lrbox}{\TitleBox}\begin{minipage}[t]{18.2cm}\raggedright\sloppy\color{Ink}\fontsize{20.7}{24}\selectfont \bfseries\color{Navy}Kriging ordinario: vincolo e modello di covarianza\end{minipage}\end{lrbox}
\begin{lrbox}{\BodyBox}\begin{minipage}[t]{10.25cm}\raggedright\sloppy\color{Ink}\fontsize{12.4}{16.4}\selectfont Il kriging ordinario cerca una previsione lineare a varianza d’errore minima, imponendo la non distorsione rispetto a una media ignota costante. I pesi derivano dal modello di covarianza e dalla geometria del campionamento; possono assumere valori negativi. Il sistema affianca alla matrice C delle covarianze il vincolo di somma unitaria, mentre c\textsubscript{0} raccoglie le covarianze fra siti osservati e posizione richiesta. La varianza di kriging è condizionata al modello specificato e richiede attenzione all’incertezza dei suoi parametri.\par\vspace{.28cm}{\centering\fontsize{13}{17}\selectfont\(\displaystyle\begin{gathered}\begin{pmatrix}C&\mathbf1\\\mathbf1^\top&0\end{pmatrix}\begin{pmatrix}\lambda\\\nu\end{pmatrix}=\begin{pmatrix}c_0\\1\end{pmatrix}\end{gathered}\)\par}\end{minipage}\end{lrbox}
\begin{lrbox}{\SideBox}\begin{minipage}[t]{7.45cm}\raggedright\sloppy\color{Ink}\fontsize{12.2}{16.1}\selectfont $\lambda$: pesi; $\mathbf1$: vettore di n elementi uguali a 1; $\nu$: moltiplicatore del vincolo. La somma dei pesi è uno. Nel reticolo con gradiente occorre prima modellare la media, eventualmente mediante kriging universale. Il laboratorio principale conserva IDW per rendere espliciti i pesi.\end{minipage}\end{lrbox}
\begin{lrbox}{\RecallBox}\begin{minipage}[t]{16.45cm}\raggedright\sloppy\color{Ink}\fontsize{10.1}{13.2}\selectfont La previsione geostatistica dipende dalla plausibilità del modello della media e della covarianza.\end{minipage}\end{lrbox}
\MeasureBox{28-TitleBox}{\TitleBox}
\MeasureBox{28-BodyBox}{\BodyBox}
\MeasureBox{28-SideBox}{\SideBox}
\MeasureBox{28-RecallBox}{\RecallBox}
\begin{frame}[plain]\hypertarget{slide28}{}
\bookmark[level=1,page=28]{Kriging ordinario: vincolo e modello di covarianza}
\begin{tikzpicture}[remember picture,overlay]\node[anchor=north east,inner sep=0] at ([xshift=19.05cm,yshift=-0.16cm]current page.north west){\color{Classical}\fontsize{8}{10}\selectfont\hyperlink{slide2}{Indice}};
\node[anchor=north west,inner sep=0] at ([xshift=0.95cm,yshift=-0.62cm]current page.north west){\usebox{\TitleBox}};
\draw[Classical,line width=1.3pt] ([xshift=.95cm,yshift=-1.68cm]current page.north west)--([xshift=2.25cm,yshift=-1.68cm]current page.north west);\node[anchor=north west,inner sep=0] at ([xshift=0.95cm,yshift=-2.12cm]current page.north west){\usebox{\BodyBox}};
\node[anchor=north west,inner sep=0] at ([xshift=11.6cm,yshift=-2.1cm]current page.north west){\color{Quantum}\bfseries\fontsize{11.5}{14}\selectfont Sviluppo e interpretazione};
\node[anchor=north west,inner sep=0] at ([xshift=11.6cm,yshift=-2.75cm]current page.north west){\usebox{\SideBox}};
\draw[Rule,line width=.5pt] ([xshift=.95cm,yshift=-10.43cm]current page.north west)--([xshift=19.05cm,yshift=-10.43cm]current page.north west);\node[anchor=north west,inner sep=0] at ([xshift=0.95cm,yshift=-10.66cm]current page.north west){\color{Classical}\bfseries\fontsize{9.4}{12}\selectfont Richiami};
\node[anchor=north west,inner sep=0] at ([xshift=2.6cm,yshift=-10.64cm]current page.north west){\usebox{\RecallBox}};
\draw[Rule,line width=.35pt] ([xshift=.95cm,yshift=-11.92cm]current page.north west)--([xshift=19.05cm,yshift=-11.92cm]current page.north west);\node[anchor=north west,inner sep=0] at ([xshift=0.95cm,yshift=-12.06cm]current page.north west){\color{Muted}\fontsize{7.7}{9.3}\selectfont M11 · Statistica spaziale};
\node[anchor=north east,inner sep=0] at ([xshift=17.65cm,yshift=-12.06cm]current page.north west){\color{Classical}\fontsize{7.3}{9.1}\selectfont Spatial Data Science, cap. 12};
\node[anchor=north east,inner sep=0] at ([xshift=19.05cm,yshift=-12.06cm]current page.north west){\color{Muted}\fontsize{7.7}{9.3}\selectfont 28/32};
\end{tikzpicture}\end{frame}
\begin{lrbox}{\TitleBox}\begin{minipage}[t]{18.2cm}\raggedright\sloppy\color{Ink}\fontsize{20.7}{24}\selectfont \bfseries\color{Navy}Separazione geografica di addestramento e test\end{minipage}\end{lrbox}
\begin{lrbox}{\BodyBox}\begin{minipage}[t]{10.25cm}\raggedright\sloppy\color{Ink}\fontsize{12.4}{16.4}\selectfont La verifica lascia fuori i ventiquattro siti con x maggiore o uguale a 1250 metri, utilizzando i quaranta siti occidentali per costruire le previsioni. Il settore di test si trova oltre l’estensione dell’addestramento lungo x, per cui l’esperimento valuta un trasferimento geografico più impegnativo dell’interpolazione fra punti circondati da misure. Una suddivisione casuale mescolerebbe invece siti prossimi, fornendo una valutazione pertinente a un’altra disponibilità di informazione. La geometria del test deve riprodurre l’impiego previsto del metodo.\par\vspace{.28cm}{\centering\fontsize{13}{17}\selectfont\(\displaystyle\begin{gathered}\mathcal T=\{i:x_i<1250\},\qquad\mathcal V=\{i:x_i\ge1250\}\end{gathered}\)\par}\end{minipage}\end{lrbox}
\begin{lrbox}{\SideBox}\begin{minipage}[t]{7.45cm}\raggedright\sloppy\color{Ink}\fontsize{12.2}{16.1}\selectfont \input{tikz/M11/D09.tex}\end{minipage}\end{lrbox}
\begin{lrbox}{\RecallBox}\begin{minipage}[t]{16.45cm}\raggedright\sloppy\color{Ink}\fontsize{10.1}{13.2}\selectfont La separazione geografica va progettata insieme alla domanda di trasferibilità.\end{minipage}\end{lrbox}
\MeasureBox{29-TitleBox}{\TitleBox}
\MeasureBox{29-BodyBox}{\BodyBox}
\MeasureBox{29-SideBox}{\SideBox}
\MeasureBox{29-RecallBox}{\RecallBox}
\begin{frame}[plain]\hypertarget{slide29}{}
\bookmark[level=1,page=29]{Separazione geografica di addestramento e test}
\begin{tikzpicture}[remember picture,overlay]\node[anchor=north east,inner sep=0] at ([xshift=19.05cm,yshift=-0.16cm]current page.north west){\color{Classical}\fontsize{8}{10}\selectfont\hyperlink{slide2}{Indice}};
\node[anchor=north west,inner sep=0] at ([xshift=0.95cm,yshift=-0.62cm]current page.north west){\usebox{\TitleBox}};
\draw[Classical,line width=1.3pt] ([xshift=.95cm,yshift=-1.68cm]current page.north west)--([xshift=2.25cm,yshift=-1.68cm]current page.north west);\node[anchor=north west,inner sep=0] at ([xshift=0.95cm,yshift=-2.12cm]current page.north west){\usebox{\BodyBox}};
\node[anchor=north west,inner sep=0] at ([xshift=11.6cm,yshift=-2.1cm]current page.north west){\color{Quantum}\bfseries\fontsize{11.5}{14}\selectfont Lettura del diagramma};
\node[anchor=north west,inner sep=0] at ([xshift=11.6cm,yshift=-2.75cm]current page.north west){\usebox{\SideBox}};
\draw[Rule,line width=.5pt] ([xshift=.95cm,yshift=-10.43cm]current page.north west)--([xshift=19.05cm,yshift=-10.43cm]current page.north west);\node[anchor=north west,inner sep=0] at ([xshift=0.95cm,yshift=-10.66cm]current page.north west){\color{Classical}\bfseries\fontsize{9.4}{12}\selectfont Richiami};
\node[anchor=north west,inner sep=0] at ([xshift=2.6cm,yshift=-10.64cm]current page.north west){\usebox{\RecallBox}};
\draw[Rule,line width=.35pt] ([xshift=.95cm,yshift=-11.92cm]current page.north west)--([xshift=19.05cm,yshift=-11.92cm]current page.north west);\node[anchor=north west,inner sep=0] at ([xshift=0.95cm,yshift=-12.06cm]current page.north west){\color{Muted}\fontsize{7.7}{9.3}\selectfont M11 · Statistica spaziale};
\node[anchor=north east,inner sep=0] at ([xshift=17.65cm,yshift=-12.06cm]current page.north west){\color{Classical}\fontsize{7.3}{9.1}\selectfont };
\node[anchor=north east,inner sep=0] at ([xshift=19.05cm,yshift=-12.06cm]current page.north west){\color{Muted}\fontsize{7.7}{9.3}\selectfont 29/32};
\end{tikzpicture}\end{frame}
\begin{lrbox}{\TitleBox}\begin{minipage}[t]{18.2cm}\raggedright\sloppy\color{Ink}\fontsize{20.7}{24}\selectfont \bfseries\color{Navy}Grandezza, segno e distribuzione degli errori\end{minipage}\end{lrbox}
\begin{lrbox}{\BodyBox}\begin{minipage}[t]{10.25cm}\raggedright\sloppy\color{Ink}\fontsize{12.4}{16.4}\selectfont Gli errori, definiti come osservato meno previsto, distinguono sottostime e sovrastime. MAE e RMSE riassumono la loro grandezza, mentre la media degli errori descrive la direzione prevalente. La carta permette di individuare una struttura geografica che tali sintesi potrebbero nascondere. Il risultato riguarda il settore Est e la potenza IDW prefissata; modifiche suggerite dal test richiederebbero una nuova verifica indipendente.\par\vspace{.28cm}{\centering\fontsize{13}{17}\selectfont\(\displaystyle\begin{gathered}e_i=z_i-\widehat z_i,\quad\operatorname{MAE}=\frac1m\sum_i|e_i|\\\operatorname{RMSE}=\sqrt{\frac1m\sum_ie_i^2}\end{gathered}\)\par}\end{minipage}\end{lrbox}
\begin{lrbox}{\SideBox}\begin{minipage}[t]{7.45cm}\raggedright\sloppy\color{Ink}\fontsize{12.2}{16.1}\selectfont \includegraphics[width=7.45cm]{figure/M11/F09.pdf}\par\vspace{.16cm}{\fontsize{9.5}{12.6}\selectfont m = 24 siti di test; errori in u.a. Grigio: 40 siti di stima; tratteggio: separazione geografica.}\end{minipage}\end{lrbox}
\begin{lrbox}{\RecallBox}\begin{minipage}[t]{16.45cm}\raggedright\sloppy\color{Ink}\fontsize{10.1}{13.2}\selectfont La carta degli errori integra le metriche aggregate e rende visibile il dominio della verifica.\end{minipage}\end{lrbox}
\MeasureBox{30-TitleBox}{\TitleBox}
\MeasureBox{30-BodyBox}{\BodyBox}
\MeasureBox{30-SideBox}{\SideBox}
\MeasureBox{30-RecallBox}{\RecallBox}
\begin{frame}[plain]\hypertarget{slide30}{}
\bookmark[level=1,page=30]{Grandezza, segno e distribuzione degli errori}
\begin{tikzpicture}[remember picture,overlay]\node[anchor=north east,inner sep=0] at ([xshift=19.05cm,yshift=-0.16cm]current page.north west){\color{Classical}\fontsize{8}{10}\selectfont\hyperlink{slide2}{Indice}};
\node[anchor=north west,inner sep=0] at ([xshift=0.95cm,yshift=-0.62cm]current page.north west){\usebox{\TitleBox}};
\draw[Classical,line width=1.3pt] ([xshift=.95cm,yshift=-1.68cm]current page.north west)--([xshift=2.25cm,yshift=-1.68cm]current page.north west);\node[anchor=north west,inner sep=0] at ([xshift=0.95cm,yshift=-2.12cm]current page.north west){\usebox{\BodyBox}};
\node[anchor=north west,inner sep=0] at ([xshift=11.6cm,yshift=-2.1cm]current page.north west){\color{Quantum}\bfseries\fontsize{11.5}{14}\selectfont Evidenza grafica};
\node[anchor=north west,inner sep=0] at ([xshift=11.6cm,yshift=-2.75cm]current page.north west){\usebox{\SideBox}};
\draw[Rule,line width=.5pt] ([xshift=.95cm,yshift=-10.43cm]current page.north west)--([xshift=19.05cm,yshift=-10.43cm]current page.north west);\node[anchor=north west,inner sep=0] at ([xshift=0.95cm,yshift=-10.66cm]current page.north west){\color{Classical}\bfseries\fontsize{9.4}{12}\selectfont Richiami};
\node[anchor=north west,inner sep=0] at ([xshift=2.6cm,yshift=-10.64cm]current page.north west){\usebox{\RecallBox}};
\draw[Rule,line width=.35pt] ([xshift=.95cm,yshift=-11.92cm]current page.north west)--([xshift=19.05cm,yshift=-11.92cm]current page.north west);\node[anchor=north west,inner sep=0] at ([xshift=0.95cm,yshift=-12.06cm]current page.north west){\color{Muted}\fontsize{7.7}{9.3}\selectfont M11 · Statistica spaziale};
\node[anchor=north east,inner sep=0] at ([xshift=17.65cm,yshift=-12.06cm]current page.north west){\color{Classical}\fontsize{7.3}{9.1}\selectfont };
\node[anchor=north east,inner sep=0] at ([xshift=19.05cm,yshift=-12.06cm]current page.north west){\color{Muted}\fontsize{7.7}{9.3}\selectfont 30/32};
\end{tikzpicture}\end{frame}
\begin{lrbox}{\TitleBox}\begin{minipage}[t]{18.2cm}\raggedright\sloppy\color{Ink}\fontsize{20.7}{24}\selectfont \bfseries\color{Navy}Esportazione e controllo del contenuto geografico\end{minipage}\end{lrbox}
\begin{lrbox}{\BodyBox}\begin{minipage}[t]{10.25cm}\raggedright\sloppy\color{Ink}\fontsize{12.4}{16.4}\selectfont La consegna comprende dati, sistema di riferimento, identificativi, codice e parametri, insieme alle carte delle osservazioni, delle previsioni e degli errori. Il GeoJSON viene riletto per verificare che geometrie e attributi utili siano stati conservati. La corrispondenza deve essere controllata per identificativo quando l’ordine delle righe può cambiare; nel frammento l’ordine resta quello della scrittura. Registrare le versioni di R, sf, GDAL e PROJ permette inoltre di ricostruire il contesto delle operazioni geometriche.\end{minipage}\end{lrbox}
\begin{lrbox}{\SideBox}\begin{minipage}[t]{7.45cm}\raggedright\sloppy\color{Ink}\fontsize{12.2}{16.1}\selectfont \par\vspace{.25cm}\lstinputlisting[language=R]{estratti/M11/geojson.R}\end{minipage}\end{lrbox}
\begin{lrbox}{\RecallBox}\begin{minipage}[t]{16.45cm}\raggedright\sloppy\color{Ink}\fontsize{10.1}{13.2}\selectfont Una consegna geografica è verificabile quando preserva dati, riferimento e procedimento di analisi.\end{minipage}\end{lrbox}
\MeasureBox{31-TitleBox}{\TitleBox}
\MeasureBox{31-BodyBox}{\BodyBox}
\MeasureBox{31-SideBox}{\SideBox}
\MeasureBox{31-RecallBox}{\RecallBox}
\begin{frame}[plain]\hypertarget{slide31}{}
\bookmark[level=0,page=31]{Consegna e studio autonomo}
\bookmark[level=1,page=31]{Esportazione e controllo del contenuto geografico}
\begin{tikzpicture}[remember picture,overlay]\node[anchor=north east,inner sep=0] at ([xshift=19.05cm,yshift=-0.16cm]current page.north west){\color{Classical}\fontsize{8}{10}\selectfont\hyperlink{slide2}{Indice}};
\node[anchor=north west,inner sep=0] at ([xshift=0.95cm,yshift=-0.62cm]current page.north west){\usebox{\TitleBox}};
\draw[Classical,line width=1.3pt] ([xshift=.95cm,yshift=-1.68cm]current page.north west)--([xshift=2.25cm,yshift=-1.68cm]current page.north west);\node[anchor=north west,inner sep=0] at ([xshift=0.95cm,yshift=-2.12cm]current page.north west){\usebox{\BodyBox}};
\node[anchor=north west,inner sep=0] at ([xshift=11.6cm,yshift=-2.1cm]current page.north west){\color{Quantum}\bfseries\fontsize{11.5}{14}\selectfont Implementazione in R};
\node[anchor=north west,inner sep=0] at ([xshift=11.6cm,yshift=-2.75cm]current page.north west){\usebox{\SideBox}};
\draw[Rule,line width=.5pt] ([xshift=.95cm,yshift=-10.43cm]current page.north west)--([xshift=19.05cm,yshift=-10.43cm]current page.north west);\node[anchor=north west,inner sep=0] at ([xshift=0.95cm,yshift=-10.66cm]current page.north west){\color{Classical}\bfseries\fontsize{9.4}{12}\selectfont Richiami};
\node[anchor=north west,inner sep=0] at ([xshift=2.6cm,yshift=-10.64cm]current page.north west){\usebox{\RecallBox}};
\draw[Rule,line width=.35pt] ([xshift=.95cm,yshift=-11.92cm]current page.north west)--([xshift=19.05cm,yshift=-11.92cm]current page.north west);\node[anchor=north west,inner sep=0] at ([xshift=0.95cm,yshift=-12.06cm]current page.north west){\color{Muted}\fontsize{7.7}{9.3}\selectfont M11 · Statistica spaziale};
\node[anchor=north east,inner sep=0] at ([xshift=17.65cm,yshift=-12.06cm]current page.north west){\color{Classical}\fontsize{7.3}{9.1}\selectfont };
\node[anchor=north east,inner sep=0] at ([xshift=19.05cm,yshift=-12.06cm]current page.north west){\color{Muted}\fontsize{7.7}{9.3}\selectfont 31/32};
\end{tikzpicture}\end{frame}
\begin{lrbox}{\TitleBox}\begin{minipage}[t]{18.2cm}\raggedright\sloppy\color{Ink}\fontsize{20.7}{24}\selectfont \bfseries\color{Navy}Percorso di laboratorio e problema di trasferimento\end{minipage}\end{lrbox}
\begin{lrbox}{\BodyBox}\begin{minipage}[t]{10.25cm}\raggedright\sloppy\color{Ink}\fontsize{12.4}{16.4}\selectfont I laboratori N01–N07 affrontano coordinate, distanze e rappresentazione; N08–N12 definiscono appartenenze, buffer e pesi; N13–N15 trattano Moran, permutazioni e semivarianza; N16–N20 completano copertura, aggregazione, interpolazione, verifica ed esportazione. Per lo studio autonomo si consideri una concentrazione misurata soltanto presso una strada, che debba essere prevista su un versante remoto. L’argomentazione deve precisare popolazione, supporto, accessibilità, differenze ambientali e geometria della validazione, prima di scegliere un interpolatore.\end{minipage}\end{lrbox}
\begin{lrbox}{\SideBox}\begin{minipage}[t]{7.45cm}\raggedright\sloppy\color{Ink}\fontsize{12.2}{16.1}\selectfont \input{tikz/M11/D10.tex}\end{minipage}\end{lrbox}
\begin{lrbox}{\RecallBox}\begin{minipage}[t]{16.45cm}\raggedright\sloppy\color{Ink}\fontsize{10.1}{13.2}\selectfont Il trasferimento a un territorio diverso richiede ipotesi sulle differenze fra dominio osservato e dominio previsto.\end{minipage}\end{lrbox}
\MeasureBox{32-TitleBox}{\TitleBox}
\MeasureBox{32-BodyBox}{\BodyBox}
\MeasureBox{32-SideBox}{\SideBox}
\MeasureBox{32-RecallBox}{\RecallBox}
\begin{frame}[plain]\hypertarget{slide32}{}
\bookmark[level=1,page=32]{Percorso di laboratorio e problema di trasferimento}
\begin{tikzpicture}[remember picture,overlay]\node[anchor=north east,inner sep=0] at ([xshift=19.05cm,yshift=-0.16cm]current page.north west){\color{Classical}\fontsize{8}{10}\selectfont\hyperlink{slide2}{Indice}};
\node[anchor=north west,inner sep=0] at ([xshift=0.95cm,yshift=-0.62cm]current page.north west){\usebox{\TitleBox}};
\draw[Classical,line width=1.3pt] ([xshift=.95cm,yshift=-1.68cm]current page.north west)--([xshift=2.25cm,yshift=-1.68cm]current page.north west);\node[anchor=north west,inner sep=0] at ([xshift=0.95cm,yshift=-2.12cm]current page.north west){\usebox{\BodyBox}};
\node[anchor=north west,inner sep=0] at ([xshift=11.6cm,yshift=-2.1cm]current page.north west){\color{Quantum}\bfseries\fontsize{11.5}{14}\selectfont Lettura del diagramma};
\node[anchor=north west,inner sep=0] at ([xshift=11.6cm,yshift=-2.75cm]current page.north west){\usebox{\SideBox}};
\draw[Rule,line width=.5pt] ([xshift=.95cm,yshift=-10.43cm]current page.north west)--([xshift=19.05cm,yshift=-10.43cm]current page.north west);\node[anchor=north west,inner sep=0] at ([xshift=0.95cm,yshift=-10.66cm]current page.north west){\color{Classical}\bfseries\fontsize{9.4}{12}\selectfont Richiami};
\node[anchor=north west,inner sep=0] at ([xshift=2.6cm,yshift=-10.64cm]current page.north west){\usebox{\RecallBox}};
\draw[Rule,line width=.35pt] ([xshift=.95cm,yshift=-11.92cm]current page.north west)--([xshift=19.05cm,yshift=-11.92cm]current page.north west);\node[anchor=north west,inner sep=0] at ([xshift=0.95cm,yshift=-12.06cm]current page.north west){\color{Muted}\fontsize{7.7}{9.3}\selectfont M11 · Statistica spaziale};
\node[anchor=north east,inner sep=0] at ([xshift=17.65cm,yshift=-12.06cm]current page.north west){\color{Classical}\fontsize{7.3}{9.1}\selectfont };
\node[anchor=north east,inner sep=0] at ([xshift=19.05cm,yshift=-12.06cm]current page.north west){\color{Muted}\fontsize{7.7}{9.3}\selectfont 32/32};
\end{tikzpicture}\end{frame}
\end{document}